%% relatorio_2809.tex
%% Relatorio de investigación (28/09/2026):
%%   Cap. 2  Revisión de literatura sobre SAR biestático y ángulo de Brewster
%%   Cap. 3-6 Análisis energético monoestático vs. biestático especular de
%%            objetos tridimensionales enterrados (metodología, resultados,
%%            discusión y conclusiones)
%% Compilar: pdflatex relatorio_2809.tex (dos pasadas)

\documentclass[11pt,a4paper]{report}

% ─── paquetes ────────────────────────────────────────────────────────────────
\usepackage[utf8]{inputenc}
\usepackage[T1]{fontenc}
\usepackage[spanish,es-tabla,es-noshorthands]{babel}
\usepackage{lmodern}
\usepackage{geometry}
\geometry{left=2.7cm, right=2.5cm, top=2.7cm, bottom=2.7cm}

\usepackage{amsmath,amssymb}
\usepackage{bm}
\usepackage{booktabs}
\usepackage{array}
\usepackage{longtable}
\usepackage{multirow}
\usepackage{enumitem}
\usepackage{xcolor}
\usepackage{graphicx}
\usepackage{caption}
\usepackage{subcaption}
\usepackage{float}
\usepackage[colorlinks=true,linkcolor=blue!50!black,citecolor=blue!50!black,urlcolor=blue!50!black]{hyperref}
\usepackage{titlesec}

\graphicspath{{figures/relatorio_2809/}}
\captionsetup{font=small, labelfont=bf}

\titleformat{\chapter}[display]{\normalfont\huge\bfseries}{\chaptertitlename\ \thechapter}{12pt}{\Huge}
\titlespacing*{\chapter}{0pt}{-20pt}{25pt}
\titleformat{\section}{\normalfont\Large\bfseries}{\thesection}{1em}{}
\titleformat{\subsection}{\normalfont\large\bfseries}{\thesubsection}{1em}{}

\setlength{\parindent}{0pt}
\setlength{\emergencystretch}{3em}
\setlength{\parskip}{0.5\baselineskip}

\newcommand{\dB}{\,\mathrm{dB}}
\newcommand{\dBm}{\,\mathrm{dBm}}
\newcommand{\dBsm}{\,\mathrm{dBsm}}
\newcommand{\vect}[1]{\bm{#1}}
\newcommand{\uv}[1]{\hat{\bm{#1}}}
\newcommand{\sig}{$^{\ast}$}
\newcolumntype{L}[1]{>{\raggedright\arraybackslash}p{#1}}
\newcolumntype{C}[1]{>{\centering\arraybackslash}p{#1}}
\newcommand{\si}{\checkmark}
\newcommand{\pa}{$\circ$}
\newcommand{\no}{--}
\newcommand{\grupo}[1]{\addlinespace[3pt]\multicolumn{6}{@{}l}{\textit{#1}}\\*\addlinespace[1pt]}

\title{\bfseries SAR biestático para sensado subsuperficial:\\[4pt]
revisión de literatura y análisis energético monoestático\\
frente a biestático especular con el ángulo de Brewster\\[10pt]
\large Relatorio de investigación}
\author{David Atencia}
\date{28 de septiembre de 2026}

\begin{document}
\maketitle

% ══════════════════════════════════════════════════════════════════════════
\begin{abstract}
\noindent
Este relatorio reúne dos contribuciones de la investigación doctoral sobre
radar de apertura sintética (SAR) biestático para la detección de blancos
enterrados. En primer lugar, presenta una revisión de 22 trabajos de la
literatura sobre SAR biestático ---fundamentos geométricos, formación de
imagen, tomografía 3D, sistemas aerotransportados y explotación del ángulo
de Brewster---, que concluye que ningún trabajo combina SAR biestático
tomográfico, refracción en la interfaz aire--suelo y optimización de la
geometría de adquisición guiada por el ángulo de Brewster. En segundo lugar,
cuantifica mediante simulación el balance energético de dos configuraciones
de adquisición sobre un arco de radio constante de 100\,m: monoestática y
biestática especular (receptor en posición simétrica al transmisor respecto
de la vertical del blanco). Se modela la propagación con refracción de Snell
y transmitancia de Fresnel en polarización TM, y la sección eficaz radar
(RCS) biestática de cinco sólidos conductores (cubo, caja, cilindro, esfera y
cono) mediante Óptica Física evaluada en el medio, a 425\,MHz y con índice
de refracción del suelo $n=2$. Con el blanco en orientación nominal, el modo
especular recibe entre $+6$ y $+10\dB$ más energía que el monoestático en los
objetos con una cara plana horizontal, debido a que la bisectriz
transmisor--receptor permanece vertical y mantiene la reflexión especular de
esa cara. Sin embargo, un experimento de Monte Carlo con $10^4$ orientaciones
aleatorias e independientes por posición muestra que, promediando sobre la
orientación del blanco, la diferencia energética entre ambos modos es
despreciable para todos los objetos (entre $-0.37$ y $+0.21\dB$). El
ángulo de Brewster aporta una ganancia de transmitancia de solo $+1.02\dB$
respecto del nadir; su principal beneficio en la configuración especular es
la anulación de la reflexión especular de la superficie, que mejora la
relación señal-reflexión de superficie en más de $40\dB$ en el entorno
inmediato del ángulo de Brewster.
\end{abstract}

\tableofcontents

% ══════════════════════════════════════════════════════════════════════════
\chapter{Introducción}
\label{cap:intro}

El radar de apertura sintética (SAR) aplicado a la detección de blancos
enterrados debe superar dos barreras físicas que no existen en la
observación de superficie: la pérdida de potencia en la interfaz aire--suelo
y la reflexión de la propia superficie, que compite con el eco del blanco.
En una geometría biestática, en la que transmisor (Tx) y receptor (Rx) se
encuentran en plataformas distintas ---por ejemplo, dos vehículos aéreos no
tripulados---, la libertad adicional de posicionar el receptor abre la
posibilidad de elegir geometrías que favorezcan la transmisión de energía
hacia el subsuelo y su retorno. El ángulo de Brewster, para el que la
reflexión de una onda con polarización TM en una interfaz dieléctrica se
anula y la transmisión es total, es el candidato físico natural para guiar
esa elección.

Este relatorio aborda dos preguntas complementarias:

\begin{enumerate}[itemsep=2pt]
  \item \textbf{¿Qué se ha publicado?} ¿Existe en la literatura algún
  trabajo que combine SAR biestático tomográfico, propagación con
  refracción en la interfaz aire--suelo y optimización de la geometría de
  adquisición guiada por el ángulo de Brewster? (Capítulo~\ref{cap:revision}.)
  \item \textbf{¿Cuánta energía se gana realmente?} Para blancos
  tridimensionales enterrados, ¿recibe más energía una configuración
  biestática especular que una monoestática cuando ambas recorren la misma
  trayectoria angular? ¿Qué parte de la diferencia se debe al ángulo de
  Brewster y qué parte a la dispersión del propio blanco? ¿Se mantiene la
  conclusión cuando la orientación del blanco es desconocida?
  (Capítulos~\ref{cap:metodologia} a~\ref{cap:discusion}.)
\end{enumerate}

La segunda pregunta es relevante para el diseño de trayectorias: si la
ventaja energética del modo biestático dependiera críticamente de la
orientación del blanco, una estrategia de vuelo optimizada para un único
supuesto geométrico podría no ser robusta en campo.

\textbf{Organización.} El Capítulo~\ref{cap:revision} presenta la revisión
de literatura. El Capítulo~\ref{cap:metodologia} describe el escenario, los
modelos de propagación, de ecuación radar y de sección eficaz, y el diseño
del experimento de Monte Carlo. El Capítulo~\ref{cap:resultados} presenta
los resultados de validación, de orientación nominal, de los diagramas de
dispersión 3D y del Monte Carlo. El Capítulo~\ref{cap:discusion} interpreta
físicamente los resultados, los sitúa respecto de la literatura y discute
sus limitaciones. El Capítulo~\ref{cap:conclusiones} resume las conclusiones
y el trabajo futuro. El Apéndice~\ref{ap:figuras} reúne las figuras de todos
los objetos simulados.

% ══════════════════════════════════════════════════════════════════════════
\chapter{Revisión de literatura: sistemas SAR biestáticos}
\label{cap:revision}

\section{Introducción y criterio de selección}
\label{sec:rev_intro}

Este capítulo reúne 22 trabajos de la literatura que abordan, en distintos
niveles de la cadena de procesado SAR (geometría, resolución, formación de
imagen, tomografía 3D, sistemas experimentales y explotación del ángulo de
Brewster), el caso en que el transmisor y el receptor \textbf{no comparten
la misma plataforma}. Esta condición es \emph{imprescindible} para la
selección de los bloques A--D, ya que define el marco de la presente
investigación: SAR biestático helicoidal con propagación en uno o varios
medios (aire--suelo). El bloque E amplía deliberadamente la búsqueda al caso
monoestático, por tratarse de una comprobación de novedad específica
(véase la Sección~\ref{sec:brewster}).

Se priorizaron publicaciones recientes (2018--2025) para capturar el estado
del arte actual, complementadas con tres referencias fundacionales
(Arikan~\&~Munson 1989, Ding~\&~Munson 2002, Horne~\&~Yates 2002) que forman
parte del marco teórico de la tesis y se incluyen para contextualizar la
evolución del campo. Cada entrada presenta una síntesis del resumen original
y el beneficio concreto de su lectura o incorporación para la presente
investigación.

Los trabajos se agrupan en cinco bloques temáticos:
\begin{enumerate}[itemsep=2pt]
  \item Fundamentos geométricos y resolución espacial biestática (Sección~\ref{sec:fundamentos}).
  \item Formación de imagen y algoritmos de reconstrucción biestática (Sección~\ref{sec:algoritmos}).
  \item Tomografía SAR biestática y multiestática 3D (Sección~\ref{sec:tomografia}).
  \item Sistemas biestáticos aerotransportados y UAV (Sección~\ref{sec:uav}).
  \item Ángulo de Brewster y optimización de potencia recibida, biestático y monoestático (Sección~\ref{sec:brewster}).
\end{enumerate}

% ──────────────────────────────────────────────────────────────────────────
\section{Fundamentos geométricos y resolución espacial biestática}
\label{sec:fundamentos}

\begin{enumerate}[label=\textbf{[A\arabic*]}, leftmargin=2.2cm, itemsep=10pt]

\item Horne, A.M.; Yates, G. (2002). \textit{Bistatic synthetic aperture
radar.} En \emph{Proc. RADAR 2002}, Edimburgo, IEE Publication No.\,490,
pp.\,6--10. \cite{horne2002}\\
\textbf{Síntesis del resumen:} el SAR es cada vez más importante en
vigilancia terrestre militar por su capacidad de operar en cualquier clima,
de día y de noche, y de detectar, clasificar y geolocalizar objetos a largo
alcance. El biestático (Tx/Rx en plataformas separadas) se ve como un medio
de contrarrestar la vulnerabilidad electrónica, manteniendo la plataforma
receptora encubierta mientras el transmisor ---más vulnerable--- puede
desplegarse más lejos. El artículo repasa la investigación de QinetiQ sobre
las técnicas de procesado para SAR biestático, los problemas fundamentales
que introduce (sincronización, ruido de fase, ejes de imagen no ortogonales)
y cómo superarlos.\\
\textbf{Beneficio:} proporciona el marco conceptual completo que sostiene el
modelo geométrico biestático de esta investigación; ayuda a anticipar
problemas prácticos (sincronización Tx/Rx, ruido de fase acumulado) antes de
un eventual despliegue real con dos UAV.

\item Homer, J.; Donskoi, E.; Mojarrabi, B.; Palmer, J.; Kubik, K. (2005).
\textit{Three-dimensional bistatic synthetic aperture radar imaging system:
spatial resolution analysis.} \emph{IEE Proceedings -- Radar, Sonar and
Navigation}, 152(6), 391--394. \cite{homer2005}\\
\textbf{Síntesis del resumen:} presentan un marco matemático para analizar
la capacidad de resolución espacial de un sistema SAR biestático 3D
general, en el que transmisor, receptor y objeto (o región de interés)
pueden ser no coplanares y en el que el transmisor y/o el receptor pueden
estar en movimiento. Con aplicaciones de observación de la Tierra en mente,
el marco se usa para caracterizar cuantitativamente la resolución 2D en el
plano del terreno del sistema de imagen 3D; los resultados teóricos se
confirman mediante simulación.\\
\textbf{Beneficio:} es la referencia teórica más cercana a la derivación de
las resoluciones horizontal y vertical, $\delta_{xy}$ y $\delta_z$, de esta
investigación, pero para geometría arbitraria (no solo circular o
helicoidal); permite generalizar el análisis de resolución a trayectorias no
ideales (deriva, viento) y validar la fórmula de resolución en el plano del
terreno con un caso 3D genérico.

\item Garza, G.; Qiao, Z. (2010). \textit{Resolution analysis of bistatic
SAR.} En \emph{Proc. SPIE 8021, Radar Sensor Technology XV}, 80210F.
\cite{garza2010}\\
\textbf{Síntesis del resumen:} analizan la resolución de imágenes SAR
biestáticas de objetos estacionarios reconstruidas mediante un método de
inversión por retroproyección filtrada, derivado de un modelo de ecuación de
onda escalar, que tiene en cuenta los patrones de haz de la antena y
trayectorias de vuelo arbitrarias. El análisis se realiza examinando la
variedad (\emph{manifold}) de recolección de datos para distintas geometrías
experimentales y parámetros del sistema.\\
\textbf{Beneficio:} su función de dispersión puntual (PSF), ligada
explícitamente al patrón de haz de la antena, permite refinar el modelo de
resolución de esta investigación incorporando el diagrama de radiación real
de las antenas Tx/Rx, que las fórmulas actuales asumen idealizado.

\item Chen, Y.; Chen, R.; Liu, H.; Guo, J.; Wang, Y.; Zhang, J. (2022).
\textit{Research on configuration constraints of airborne bistatic SARs.}
\emph{Sensors}, 22(17), 6534.
\href{https://doi.org/10.3390/s22176534}{doi:10.3390/s22176534}
\cite{chen2022}\\
\textbf{Síntesis del resumen:} a partir del análisis de la geometría de
imagen del SAR biestático aerotransportado, emplean un algoritmo de
\emph{chirp scaling} no lineal extendido para simular y verificar el efecto
de imagen. Proponen un modelo de resolución 2D biestático basado en teoría
de gradientes y analizan las restricciones de los parámetros de trayectoria
multiplataforma que satisfacen la precisión de imagen requerida. Mediante
simulación con trayectorias cooperativas en distintos escenarios, revelan la
envolvente de restricción de configuración espacial entre los vehículos
para lograr la resolución óptima.\\
\textbf{Beneficio:} su ``envolvente de configuración válida'' entre Tx y Rx
es directamente aplicable para acotar el espacio de búsqueda del
optimizador de trayectoria del receptor, hoy basado en un esquema de
arranques múltiples sin restricción geométrica derivada analíticamente, lo
que podría acelerar su convergencia.

\end{enumerate}

% ──────────────────────────────────────────────────────────────────────────
\section{Formación de imagen y algoritmos de reconstrucción biestática}
\label{sec:algoritmos}

\begin{enumerate}[label=\textbf{[B\arabic*]}, leftmargin=2.2cm, itemsep=10pt]

\item Arikan, O.; Munson, D.C.\,Jr. (1989). \textit{A tomographic
formulation of bistatic synthetic aperture radar.} En Porter, W.A.; Kak,
S.C. (eds.), \emph{Advances in Communications and Signal Processing},
Lecture Notes in Control and Information Sciences, vol.\,129, pp.\,418--431,
Springer. (Presentado en \emph{Proc. ComCon '88}, Baton Rouge, LA, oct.\,1988.)
\cite{arikan1989}\\
\textbf{Síntesis del resumen:} un radar biestático es aquel con antenas
transmisora y receptora separadas por una distancia comparable a la
distancia al objetivo; el concepto no es nuevo pero fue eclipsado durante
décadas por el monoestático. En los años previos hubo un interés militar
considerable en el SAR biestático de modo \emph{spotlight} (BSSAR), en el
que un transmisor de alta potencia opera a distancia segura mientras un
receptor encubierto vuela cerca del área objetivo. Dan una derivación
básica del BSSAR desde primeros principios, mostrando ---de forma análoga al
caso monoestático (MSSAR)--- que puede explicarse mediante el teorema de
corte-proyección de la tomografía asistida por computador, sin depender del
efecto Doppler: para cada ángulo de observación, los datos BSSAR
demodulados son muestras aproximadamente limitadas en banda de una rebanada
de la transformada de Fourier 2D de la reflectividad del objetivo. Examinan
la ubicación de las muestras en el dominio de Fourier y la curvatura del
frente de onda para varios casos especiales de movimiento de Tx y Rx.\\
\textbf{Beneficio:} es la base tomográfica (dominio del número de onda) que
sostiene el marco teórico de la tesis; su análisis de la curvatura del
frente de onda permite verificar que la aproximación de campo lejano del
modelo de propagación multicapa sigue siendo válida para las distancias
Tx--blanco del escenario con UAV.

\item Ding, Y.; Munson, D.C.\,Jr. (2002). \textit{A fast back-projection
algorithm for bistatic SAR imaging.} En \emph{Proc. IEEE Int. Conf. on Image
Processing (ICIP)}, vol.\,II, pp.\,449--452. \cite{ding2002}\\
\textbf{Síntesis del resumen:} con un modelo de campo lejano, el SAR
biestático adquiere datos de Fourier en una grilla no cartesiana poco
habitual. Los algoritmos previos se basaban principalmente en la
reconstrucción directa de Fourier (DFR) para aprovechar la FFT, pero la
cobertura irregular de los datos y la interpolación 2D en el dominio de
Fourier pueden afectar la precisión de la reconstrucción. La
retroproyección (BP) evita estos problemas pero tiene un costo
computacional enorme, $O(N^3)$ frente a $O(N^2\log_2 N)$ de la DFR.
Presentan un algoritmo BP rápido para SAR biestático, inspirado en un
algoritmo rápido de BP para tomografía, que logra un costo del mismo orden
que la DFR y que además puede usarse para imagen de campo cercano; validan
el rendimiento mediante simulación.\\
\textbf{Beneficio:} es el antecedente algorítmico directo de la
retroproyección empleada en la cadena de procesamiento de esta
investigación; su versión rápida (jerárquica) permitiría escalar la
reconstrucción a grillas volumétricas más finas sin el costo $O(N^3)$
actual.

\item Zhang, H.; Tang, J.; Wang, R.; Deng, Y.; Wang, W.; Li, N. (2018).
\textit{An accelerated backprojection algorithm for monostatic and bistatic
SAR processing.} \emph{Remote Sensing}, 10(1), 140.
\href{https://doi.org/10.3390/rs10010140}{doi:10.3390/rs10010140}
\cite{zhang2018}\\
\textbf{Síntesis del resumen:} proponen un algoritmo BP acelerado a partir
de Block\_FFBP en el que un número fijo de ``pivotes'' ---en vez de centros
de haz que crecen con cada etapa de factorización--- construye la relación
de retardo de propagación, y los datos de rango se procesan como un bloque
completo, sin partición en la dimensión de rango, gracias a un esquema
integrado de procesamiento en rango. El esquema se extiende a SAR
biestático (BiBulk\_FFBP) y monoestático (MoBulk\_FFBP), transfiriendo el
cálculo del rango oblicuo del sistema polar al cartesiano. Validan el
algoritmo con datos reales del experimento SAR biestático
espacial/estacionario TerraSAR-X en modo \emph{staring spotlight} y con un
experimento SAR aerotransportado de 2016.\\
\textbf{Beneficio:} ruta de aceleración computacional validada con datos
reales de una misión biestática espacial; referencia de eficiencia
aplicable a la retroproyección volumétrica de esta investigación si se
adopta un esquema jerárquico (FFBP).

\item Wang, Z.; Liu, M.; Ai, G.; Wang, P.; Lv, K. (2020). \textit{Focusing of
bistatic SAR with curved trajectory based on extended azimuth nonlinear
chirp scaling.} \emph{IEEE Transactions on Geoscience and Remote Sensing},
58(6), 4160--4179. \cite{wang2020}\\
\textbf{Síntesis del resumen:} proponen un modelo de rango de movimiento de
alto orden y un algoritmo de \emph{chirp scaling} no lineal acimutal
extendido (EANLCS) con corrección residual del centroide Doppler, para
enfocar datos SAR biestáticos con trayectoria curva de las plataformas. El
trabajo aborda el acoplamiento rango--acimut adicional y la variación
espacial de la migración de celda de rango y de la tasa Doppler causadas por
la trayectoria curva, efectos que los algoritmos tradicionales ---derivados
para trayectorias rectilíneas de velocidad uniforme--- no pueden procesar.\\
\textbf{Beneficio:} es la referencia más cercana al caso de trayectorias
helicoidales cónicas de esta investigación, ya que ninguno de los algoritmos
biestáticos clásicos (Arikan, Ding) asume trayectoria curva; su modelo de
rango de alto orden es candidato a mejorar la precisión de fase del
procesador si en el futuro se simplifica el cálculo exacto de la distancia
por el principio de Fermat.

\item Sun, H.; Sun, Z.; Chen, T.; Miao, Y.; Wu, J.; Yang, J. (2022).
\textit{An efficient backprojection algorithm based on wavenumber-domain
spectral splicing for monostatic and bistatic SAR configurations.}
\emph{Remote Sensing}, 14(8), 1885.
\href{https://doi.org/10.3390/rs14081885}{doi:10.3390/rs14081885}
\cite{sun2022}\\
\textbf{Síntesis del resumen:} en la FBP tradicional, cada subapertura usa
su propio sistema de coordenadas polares centrado en ella, por lo que las
imágenes de subapertura deben proyectarse a un sistema uniforme antes de
sumarse coherentemente, con gran costo computacional. Proponen usar el
mismo sistema polar para todas las subaperturas y empalmar (\emph{splice})
sus imágenes en el dominio del número de onda, sumándolas directamente tras
un sobremuestreo, lo que mejora precisión y eficiencia; el método es
aplicable a configuraciones mono y biestáticas.\\
\textbf{Beneficio:} técnica de fusión en el dominio del número de onda
directamente relacionada con el análisis de la cobertura espectral
$k_{xy}$--$k_z$ de esta investigación; útil para acelerar la formación de la
imagen volumétrica sin sacrificar la precisión de fase que determina la
resolución vertical $\delta_z$.

\item Zhao, T.; Wang, J.; Cheng, Z.; Huang, Z.; Song, J. (2025). \textit{A
subimage autofocus bistatic ground Cartesian back-projection algorithm for
passive bistatic SAR based on GEO satellites.} \emph{Remote Sensing}, 17(9),
1576. \href{https://doi.org/10.3390/rs17091576}{doi:10.3390/rs17091576}
\cite{zhao2025}\\
\textbf{Síntesis del resumen:} el SAR biestático pasivo (PBSAR) basado en
señales de satélites geoestacionarios muestra potencial para imagen de alta
resolución, pero enfrenta los desafíos de la eficiencia computacional y la
compensación de errores de fase. Proponen un algoritmo de retroproyección
cartesiana en tierra con autoenfoque de subimagen (SIAF-BGCBP), basado en la
corrección del espectro de número de onda de cada subimagen, que reduce en
un 90\,\% la complejidad de la estimación de fase al optimizar la densidad
de píxeles de subimagen, manteniendo la precisión.\\
\textbf{Beneficio:} relevante si en el futuro se explora una arquitectura
biestática pasiva de bajo costo (iluminador fijo o satelital + receptor en
UAV) como alternativa al esquema Tx/Rx activo dual; su autoenfoque por
subimagen es transferible a la corrección de fase por errores de
posicionamiento GNSS del receptor.

\end{enumerate}

% ──────────────────────────────────────────────────────────────────────────
\section{Tomografía SAR biestática y multiestática 3D}
\label{sec:tomografia}

\begin{enumerate}[label=\textbf{[C\arabic*]}, leftmargin=2.2cm, itemsep=10pt]

\item Ge, N.; Zhu, X.X. (2019). \textit{Bistatic-like differential SAR
tomography.} \emph{IEEE Transactions on Geoscience and Remote Sensing},
57(8), 5883--5893. \cite{ge2019}\\
\textbf{Síntesis del resumen:} abordan la tomografía SAR diferencial
(D-TomoSAR) en áreas urbanas con adquisiciones espaciales biestáticas o
``pursuit monostatic''. Estos interferogramas no sufren decorrelación
temporal significativa ni pantalla de fase atmosférica, por lo que son
idóneos para reconstruir la elevación. Proponen un marco de dos etapas:
primero reconstruyen la elevación de los dispersores con interferogramas de
tipo biestático y luego la incorporan como información a priori
determinística para estimar la deformación a partir de interferogramas
convencionales de paso repetido.\\
\textbf{Beneficio:} su estrategia de dos etapas es transferible para
separar la localización 3D del blanco enterrado (elevación o profundidad) de
la estimación de su reflectividad, lo que podría reducir el espacio de
búsqueda de la retroproyección volumétrica.

\item Chen, Z.; Li, Y.; Hu, C. et al. (2024). \textit{Repeat-pass
space-surface bistatic SAR tomography: accurate imaging and first
experiment.} \emph{Science China Information Sciences}, 67, 192304.
\href{https://doi.org/10.1007/s11432-024-4089-2}{doi:10.1007/s11432-024-4089-2}
\cite{chenz2024}\\
\textbf{Síntesis del resumen:} la tomografía SAR biestática
espacio--superficie (SS-BiSAR) ofrece un ángulo de observación adicional al
SAR monoestático espacial y es prometedora para recuperar deformaciones con
alta precisión en regiones locales. Presentan el primer experimento de
tomografía SS-BiSAR de paso repetido, con el satélite gemelo chino LuTan-1
en banda L, adquiriendo 11 imágenes y recuperando nubes de puntos 3D de un
área edificada con $\sim$4\,m de precisión. Proponen un método de
corrección de efemérides basado en la transformada chirp-Z para corregir la
distorsión geométrica y los errores de fase interferométrica causados por
imprecisiones orbitales.\\
\textbf{Beneficio:} primera demostración experimental de tomografía
biestática 3D real en banda L, la misma banda que una de las
configuraciones del sistema estudiado en esta investigación; su corrección
de la posición de la plataforma es análoga al problema de los errores de
posicionamiento del receptor UAV que afectarían la fase de la
retroproyección en campo.

\item Wu, J.; Yang, J.; Li, Z.; Pu, W.; Sun, Z.; An, H.; Yang, Q. (2025).
\textit{Review of bistatic synthetic aperture radar imaging methods.}
\emph{Journal of Radars}, 14(5), 1115--1141.
\href{https://doi.org/10.12000/JR25067}{doi:10.12000/JR25067}
\cite{wu2025}\\
\textbf{Síntesis del resumen:} el SAR biestático emplea plataformas de
transmisión y recepción espacialmente separadas para obtener imágenes de
alta resolución de escenas terrestres y marítimas en entornos complejos,
con ventajas de configuración flexible, capacidad de ocultamiento,
resistencia a interferencias y adquisición más completa de información del
blanco. Como el modelo de eco del BiSAR difiere sustancialmente del
monoestático, el procesado de imagen es crítico. La revisión examina los
desafíos y soluciones para las configuraciones BiSAR aerotransportada, de
plataformas de alta velocidad y maniobrabilidad, espacial heterogénea y
espacial homogénea, junto con la compensación de movimiento y la imagen de
blancos móviles, y ofrece una perspectiva de tendencias futuras.\\
\textbf{Beneficio:} es el mapa más actualizado (2025) del campo completo;
su lectura íntegra es la forma más eficiente de verificar que el estado del
arte de la tesis no omite ninguna línea activa de investigación.

\item He, J.; Xie, H.; Liu, H.; Wu, Z.; Xu, B.; Zhu, N.; Lu, Z.; Qin, P.
(2025). \textit{Multi-baseline bistatic SAR three-dimensional imaging
method based on phase error calibration combining PGA and EB-ISOA.}
\emph{Remote Sensing}, 17(3), 363.
\href{https://doi.org/10.3390/rs17030363}{doi:10.3390/rs17030363}
\cite{he2025}\\
\textbf{Síntesis del resumen:} la TomoSAR monoestática multi-línea de base
sufre decorrelación temporal en escenas como bosques; además, el
\emph{jitter} de la plataforma y los errores de medición de posición
introducen errores de fase que degradan la reconstrucción. Proponen un
método de imagen 3D biestático multi-línea de base que combina la
calibración de fase PGA (\emph{Phase Gradient Autofocus}) con un nuevo
método EB-ISOA para corregir estos errores.\\
\textbf{Beneficio:} aborda exactamente el problema de \emph{jitter} de
plataforma y error de posición degradando la fase que enfrentaría cualquier
implementación real con dos UAV en trayectorias helicoidales
independientes; su combinación PGA+EB-ISOA es un candidato concreto de
autoenfoque antes de una validación experimental.

\item Andre, D.; Krishnan, V.P.; Nolan, C. (2025). \textit{Artefact analysis
of multistatic three-dimensional SAR imaging with two linear
trajectories.} arXiv:2503.11323. \cite{andre2025}\\
\textbf{Síntesis del resumen:} con la mejora de la tecnología de radar, el
SAR multiestático es hoy una posibilidad realista, por ejemplo en
constelaciones de satélites o en sistemas de aeronaves no tripuladas, lo que
exige investigar geometrías multiestáticas útiles. Realizan un análisis
microlocal de una geometría en la que transmisor y receptor se mueven
independientemente a lo largo de dos rectas que cubren una región de
interés (ROI), muestran cómo la geometría de adquisición determina la
presencia de artefactos en la imagen y proponen una condición de
\emph{time-gating} sobre los datos que garantiza la ausencia de artefactos
en la ROI, verificada con simulaciones numéricas.\\
\textbf{Beneficio:} herramienta analítica rigurosa para predecir y evitar
artefactos antes de simular; aplicar su condición de \emph{time-gating} a la
cadena de generación de señal y procesamiento permitiría filtrar a priori
qué pares de posiciones Tx/Rx de la trayectoria helicoidal aportan energía
útil al blanco, reduciendo artefactos en la retroproyección volumétrica.

\end{enumerate}

% ──────────────────────────────────────────────────────────────────────────
\section{Sistemas biestáticos aerotransportados y UAV}
\label{sec:uav}

\begin{enumerate}[label=\textbf{[D\arabic*]}, leftmargin=2.2cm, itemsep=10pt]

\item Liu, M. et al. (2024). \textit{Unmanned airborne bistatic
interferometric synthetic aperture radar data processing method using
bi-directional synchronization chain signals.} \emph{Remote Sensing},
16(5), 769. \href{https://doi.org/10.3390/rs16050769}{doi:10.3390/rs16050769}
\cite{liu2024}\\
\textbf{Síntesis del resumen:} los UAV equipados con InSAR biestático tienen
bajo costo y alta flexibilidad, útiles para cartografía y exploración de
recursos. Debido a los desafíos de sincronización espaciotemporal y errores
de movimiento, proponen un método integral de procesamiento para InSAR
biestático en UAV de pequeño porte, que integra señales de cadena de
sincronización bidireccional, compensación de errores de sincronización
temporal y de fase, refinamiento de trayectoria, imagen InSAR biestática de
alta precisión y procesamiento interferométrico. El método mejora la
precisión interferométrica en puntos de calibración de 0.66\,m a 0.42\,m.\\
\textbf{Beneficio:} es un sistema UAV biestático real y completo (hardware y
procesado) con el mismo desafío de sincronización de dos plataformas aéreas
independientes que enfrentaría una implementación física de esta
investigación; su sincronización bidireccional es un candidato directo de
diseño de hardware para una validación experimental.

\item Grathwohl, A.; Meinecke, B.; Widmann, M.; Kanz, J.; Waldschmidt, C.
(2023). \textit{UAV-based bistatic SAR-imaging using a stationary
repeater.} \emph{IEEE Journal of Microwaves}, 3(2), 625--634.
\href{https://doi.org/10.1109/JMW.2023.3253667}{doi:10.1109/JMW.2023.3253667}
\cite{grathwohl2023}\\
\textbf{Síntesis del resumen:} el SAR se apoya en el movimiento propio del
radar y en un largo tiempo de integración coherente, iluminando el blanco
desde distintos ángulos; una configuración biestática permite adquirir
simultáneamente múltiples ángulos de observación. Proponen un SAR
biestático con un repetidor activo estacionario (BR-SAR) que actúa como
segundo observador y proporciona una vía biestática con bajo costo y
esfuerzo de hardware; describen las consideraciones de procesado y hardware
y validan el enfoque con imágenes monoestáticas, biestáticas y
multiestáticas combinadas de alta resolución.\\
\textbf{Beneficio:} demuestra que un sistema biestático con UAV puede
construirse con hardware simple (repetidor estacionario) en vez de dos
plataformas activas sincronizadas; alternativa de bajo costo para una
futura validación experimental de la trayectoria de receptor optimizada.

\item Kanz, J.; Gesell, C.; Bonfert, C.; Werbunat, D.; Grathwohl, A.;
Aguilar, J.; Vossiek, M.; Waldschmidt, C. (2025). \textit{UAV-borne digital
radar system for coherent multistatic SAR imaging.} arXiv:2507.20792.
\cite{kanz2025}\\
\textbf{Síntesis del resumen:} los avances en conversores analógico-digitales
permiten arquitecturas de radar digital que muestrean directamente la señal
de radiofrecuencia, con mayor flexibilidad de forma de onda, en particular
mediante modulación OFDM. Presentan un radar digital en UAV con formas de
onda OFDM para imagen SAR multiestática coherente en banda L: un nodo
primario transmite y adquiere datos monoestáticos, mientras nodos
secundarios solo reciben, capturando la señal reflejada por la escena y una
señal de enlace lateral directo. El sistema mantiene la coherencia en tiempo
rápido y lento, y al ser pasivos los nodos secundarios permite escalar a un
enjambre de UAV. Detallan el procesamiento completo, incluida la corrección
de errores de sincronización, y lo validan con mediciones estáticas y un
experimento SAR biestático con UAV.\\
\textbf{Beneficio:} la arquitectura OFDM con un nodo Tx activo y nodos Rx
pasivos escalables es la ruta más directa para llevar el escenario
``Tx fijo + Rx optimizado'' de esta investigación a una validación de campo;
su sincronización mediante enlace lateral directo resuelve el problema de
coherencia de fase Tx/Rx, crítico para trasladar la optimización de
trayectoria multicapa de la simulación al vuelo real.

\end{enumerate}

% ──────────────────────────────────────────────────────────────────────────
\section{Ángulo de Brewster y optimización de potencia recibida (biestático y monoestático)}
\label{sec:brewster}

Esta sección responde a una búsqueda dirigida y específica: \emph{¿existe en
la literatura algún trabajo de SAR tomográfico que optimice la trayectoria o
la geometría de adquisición en función del ángulo de Brewster para
maximizar la potencia recibida, en configuración biestática o
monoestática?} Se relaja deliberadamente el criterio ``imprescindiblemente
biestático'' de las secciones A--D para cubrir también el caso monoestático.

\textbf{Resultado de la búsqueda:} no se encontró ningún trabajo que
combine (i) formación de imagen SAR \emph{tomográfica} (reconstrucción 3D)
con (ii) optimización activa de la trayectoria o de la geometría Tx/Rx en
función del ángulo de Brewster para maximizar la potencia recibida ---ni en
biestático ni en monoestático---. Los trabajos más cercanos se dividen en dos
categorías, ninguna de las cuales cierra esa brecha:

\begin{enumerate}[label=\textbf{[E\arabic*]}, leftmargin=2.2cm, itemsep=10pt]

\item Ouchi, K.; Yang, C.-S. (2017). \textit{Direct evidence and
implications of Brewster's angle damping observed by synthetic aperture
radar images.} En \emph{Proc. IGARSS 2017.} \cite{ouchi2017}\\
\textbf{Síntesis del resumen:} reportan evidencia directa del
amortiguamiento por ángulo de Brewster y sus implicaciones en el análisis
polarimétrico, observado en imágenes SAR de estructuras de concreto como
espigones y rompeolas.\\
\textbf{Categoría:} monoestático; \emph{observación pasiva} del efecto
Brewster en imágenes ya adquiridas (sin optimización de trayectoria ni
tomografía 3D).\\
\textbf{Beneficio:} confirma que el efecto Brewster es medible y
significativo en SAR monoestático de superficie con datos reales
(TerraSAR-X); constituye evidencia empírica independiente de que el
contraste de potencia en función del ángulo de incidencia explotado en esta
investigación es un fenómeno físicamente observable.

\item Ouchi, K.; Yang, C.-S. (2018). \textit{Brewster angle damping
observed in the TerraSAR-X synthetic aperture radar images of man-made
targets.} \emph{IEEE Geoscience and Remote Sensing Letters}, 15(4),
532--536. \cite{ouchi2018}\\
\textbf{Síntesis del resumen:} versión ampliada del trabajo anterior:
analizan el amortiguamiento por ángulo de Brewster en imágenes SAR
co-polarizadas de TerraSAR-X sobre blancos artificiales (estructuras de
concreto), cuantificando el fenómeno y discutiendo sus implicaciones para la
interpretación polarimétrica de imágenes de superficie.\\
\textbf{Categoría:} monoestático; análisis polarimétrico \emph{a posteriori}
(sin optimización de trayectoria ni reconstrucción 3D).\\
\textbf{Beneficio:} aporta un modelo cuantitativo del amortiguamiento por
Brewster en datos SAR reales, útil como referencia de contraste para validar
experimentalmente la ganancia de potencia predicha por el modelo de
Fresnel/Brewster de la función de costo del optimizador de trayectoria.

\item Thirion-Lefevre, L.; Guinvarc'h, R. (2018). \textit{The double
Brewster angle effect.} \emph{Comptes Rendus Physique}, 19(1--2), 43--53.
\href{https://doi.org/10.1016/j.crhy.2018.02.003}{doi:10.1016/j.crhy.2018.02.003}
\cite{thirion2018}\\
\textbf{Síntesis del resumen:} el efecto de doble ángulo de Brewster (DBE)
extiende el concepto de ángulo de Brewster a la doble reflexión sobre dos
superficies dieléctricas ortogonales, como combinación de dos
pseudo-ángulos de Brewster que ocurren en dominios de incidencia
complementarios. Demuestran teórica y experimentalmente que el DBE puede
observarse potencialmente en cualquier imagen de radar monoestático de la
superficie terrestre.\\
\textbf{Categoría:} monoestático; fenómeno físico de doble reflexión (sin
optimización de trayectoria ni SAR tomográfico).\\
\textbf{Beneficio:} generaliza el ángulo de Brewster simple de la interfaz
plana aire--suelo al caso de superficies dieléctricas ortogonales; relevante
como línea futura si el modelo se extiende a geometrías de subsuelo no
planas (zanjas, esquinas de estructuras enterradas).

\item Orfeo, D.; Zhang, Y.; Burns, D.; Miller, J.S.; Huston, D.R.; Xia, T.
(2019). \textit{Bistatic antenna configurations for air-launched ground
penetrating radar.} \emph{Journal of Applied Remote Sensing}, 13(2),
027501. \href{https://doi.org/10.1117/1.JRS.13.027501}{doi:10.1117/1.JRS.13.027501}
\cite{orfeo2019}\\
\textbf{Síntesis del resumen:} el GPR es valioso para detectar objetos
subsuperficiales con poco o ningún contenido metálico (plásticos,
cerámicas, tuberías de concreto), pero los efectos de parámetros de la
configuración de antena como la altura y el ángulo no están bien estudiados.
Mediante simulaciones y experimentos de laboratorio evalúan el efecto del
ángulo y de la altura de antena en la sensibilidad de un GPR biestático de
lanzamiento aéreo en la búsqueda de objetos no metálicos enterrados; los
resultados orientan el desarrollo de GPR aerotransportados instalados en
UAV.\\
\textbf{Categoría:} \textbf{biestático}; optimización de ángulo y altura de
antena orientada a UAV ---el más cercano en espíritu a esta
investigación---, pero sin formación de imagen SAR ni tomográfica: es un
único par Tx/Rx fijo por medición (\emph{point sensing}), no una síntesis de
apertura con trayectoria continua, y no menciona expresamente el ángulo de
Brewster como criterio de diseño, aunque el fenómeno explorado ---mínimo de
reflexión en la interfaz según ángulo y polarización--- es el mismo.\\
\textbf{Beneficio:} es el antecedente más próximo al problema concreto de
esta investigación (GPR biestático en UAV con geometría Tx/Rx variable para
maximizar la señal transmitida al subsuelo); su validación de laboratorio es
una referencia para diseñar un experimento de campo que confirme la
ganancia de potencia predicha, y su ausencia de imagen 3D subraya la brecha
que esta investigación aborda.

\end{enumerate}

\textbf{Conclusión de la búsqueda dirigida:} el ángulo de Brewster como
criterio físico para maximizar la transmisión de potencia hacia el subsuelo
está establecido en la literatura de GPR de punto único (antena fija a un
ángulo y altura óptimos) y se ha \emph{observado} de forma pasiva en imágenes
SAR monoestáticas de superficie. Sin embargo, a septiembre de 2026 no se
encontró ningún trabajo ---biestático ni monoestático--- que integre ese
criterio en un \emph{optimizador de trayectoria} acoplado a una
\emph{formación de imagen SAR tomográfica 3D} de blancos subsuperficiales,
que es la contribución de la optimización de trayectoria propuesta en esta
tesis, con función de costo
\begin{equation}
  J = -(1-\alpha)\log_{10}P_r + \alpha\log_{10}\left(\delta_{xy}\,\delta_z\right),
\end{equation}
donde $P_r$ es la potencia recibida, $\delta_{xy}$ y $\delta_z$ las
resoluciones horizontal y vertical y $\alpha\in[0,1]$ el peso relativo.
Esto refuerza la síntesis de la Sección~\ref{sec:sintesis}: la combinación
de SAR biestático tomográfico, refracción multicapa y optimización de
trayectoria por ángulo de Brewster no tiene, hasta donde permite establecer
esta búsqueda, antecedente directo en la literatura. Los
Capítulos~\ref{cap:resultados} y~\ref{cap:discusion} precisan
cuantitativamente cuánto aporta el ángulo de Brewster a la energía recibida.

% ──────────────────────────────────────────────────────────────────────────
\section{Síntesis: posicionamiento de esta investigación respecto a la literatura}
\label{sec:sintesis}

{\small
\setlength{\tabcolsep}{4pt}
\renewcommand{\arraystretch}{1.15}
\begin{longtable}{@{}L{3.6cm}C{1.0cm}C{1.0cm}C{0.8cm}C{1.6cm}L{6.0cm}@{}}
\caption{Posicionamiento de esta investigación frente a los trabajos
revisados. Criterios: \textbf{Biest.}, geometría con Tx y Rx en plataformas
distintas; \textbf{Suelo}, propagación con refracción en la interfaz
aire--suelo; \textbf{3D}, formación de imagen tridimensional o tomográfica;
\textbf{Brewster}, uso del ángulo de Brewster como criterio de diseño.
\si: sí; \pa: parcial, implícito u observado a posteriori; \no: no.}
\label{tab:sintesis}\\
\toprule
\textbf{Trabajo} & \textbf{Biest.} & \textbf{Suelo} & \textbf{3D} & \textbf{Brewster} & \textbf{Aporte principal} \\
\midrule
\endfirsthead
\toprule
\textbf{Trabajo} & \textbf{Biest.} & \textbf{Suelo} & \textbf{3D} & \textbf{Brewster} & \textbf{Aporte principal} \\
\midrule
\endhead
\bottomrule
\endfoot
\grupo{A. Fundamentos geométricos y resolución}
Horne \& Yates 2002 \cite{horne2002} & \si & \no & \no & \no & Marco conceptual; sincronización y ruido de fase \\
Homer et al.\ 2005 \cite{homer2005} & \si & \no & \si & \no & Resolución de un sistema biestático 3D general \\
Garza \& Qiao 2010 \cite{garza2010} & \si & \no & \no & \no & Resolución con patrón de antena y trayectoria arbitraria \\
Chen et al.\ 2022 \cite{chen2022} & \si & \no & \no & \no & Restricciones de configuración entre plataformas aéreas \\
\grupo{B. Formación de imagen y reconstrucción}
Arikan \& Munson 1989 \cite{arikan1989} & \si & \no & \no & \no & Formulación tomográfica (corte-proyección) del SAR biestático \\
Ding \& Munson 2002 \cite{ding2002} & \si & \no & \no & \no & Retroproyección rápida biestática \\
Zhang et al.\ 2018 \cite{zhang2018} & \si & \no & \no & \no & FFBP acelerada, validada con TerraSAR-X biestático \\
Wang et al.\ 2020 \cite{wang2020} & \si & \no & \no & \no & Enfoque con trayectoria curva (EANLCS) \\
Sun et al.\ 2022 \cite{sun2022} & \si & \no & \no & \no & Empalme espectral en el dominio del número de onda \\
Zhao et al.\ 2025 \cite{zhao2025} & \si & \no & \no & \no & Autoenfoque para biestático pasivo con iluminador GEO \\
\grupo{C. Tomografía biestática y multiestática 3D}
Ge \& Zhu 2019 \cite{ge2019} & \pa & \no & \si & \no & TomoSAR diferencial urbana con interferogramas de tipo biestático \\
Chen Z.\ et al.\ 2024 \cite{chenz2024} & \si & \no & \si & \no & Primer experimento de tomografía espacio--superficie (banda L) \\
Wu et al.\ 2025 \cite{wu2025} & \si & \no & \pa & \no & Revisión de métodos de imagen biestática \\
He et al.\ 2025 \cite{he2025} & \si & \no & \si & \no & Calibración de fase para imagen 3D multi-línea de base \\
Andre et al.\ 2025 \cite{andre2025} & \si & \no & \si & \no & Análisis microlocal de artefactos en imagen 3D multiestática \\
\grupo{D. Sistemas aerotransportados y UAV}
Liu et al.\ 2024 \cite{liu2024} & \si & \no & \no & \no & InSAR biestático en UAV con sincronización bidireccional \\
Grathwohl et al.\ 2023 \cite{grathwohl2023} & \si & \no & \no & \no & SAR biestático con UAV y repetidor estacionario \\
Kanz et al.\ 2025 \cite{kanz2025} & \si & \no & \no & \no & Radar digital OFDM multiestático en UAV \\
\grupo{E. Ángulo de Brewster y potencia recibida}
Ouchi \& Yang 2017 \cite{ouchi2017} & \no & \no & \no & \pa & Observación del amortiguamiento de Brewster en SAR \\
Ouchi \& Yang 2018 \cite{ouchi2018} & \no & \no & \no & \pa & Cuantificación en TerraSAR-X sobre blancos artificiales \\
Thirion-Lefevre \& Guinvarc'h 2018 \cite{thirion2018} & \no & \no & \no & \pa & Efecto de doble ángulo de Brewster \\
Orfeo et al.\ 2019 \cite{orfeo2019} & \si & \si & \no & \pa & GPR biestático aerotransportado; ángulo y altura de antena para un par fijo \\
\midrule
\textbf{Esta investigación} & \si & \si & \si & \si & \textbf{SAR biestático 3D con refracción multicapa y trayectoria del receptor guiada por Brewster} \\
\end{longtable}
}

Ningún trabajo revisado combina simultáneamente: (i) geometría biestática
con trayectorias helicoidales cónicas independientes para Tx y Rx, (ii)
propagación con refracción en la interfaz aire--suelo (o multicapa) y (iii)
optimización conjunta de potencia recibida (ángulo de Brewster) y
resolución 3D mediante síntesis de apertura SAR tomográfica. La búsqueda
dirigida de la Sección~\ref{sec:brewster} confirma que, incluso relajando la
condición biestática, no existe antecedente que cierre esta brecha. Esta
combinación constituye la contribución diferencial de la presente
investigación dentro del panorama actual del SAR biestático.

% ══════════════════════════════════════════════════════════════════════════
\chapter{Metodología}
\label{cap:metodologia}

\section{Objetivo del experimento numérico}

El experimento cuantifica la potencia recibida por dos configuraciones de
adquisición que recorren la misma trayectoria angular sobre un blanco
conductor enterrado:

\begin{itemize}[itemsep=2pt]
  \item \textbf{Monoestática (M):} el receptor coincide con el transmisor.
  \item \textbf{Biestática especular (B):} el receptor ocupa la posición
  simétrica del transmisor respecto del eje vertical que pasa por el blanco.
\end{itemize}

Se busca separar tres contribuciones a la diferencia de potencia entre ambas:
(i) la transmitancia de la interfaz aire--suelo, que depende del ángulo de
incidencia y es máxima en el ángulo de Brewster; (ii) la sección eficaz
radar (RCS) del blanco, que depende de la geometría de dispersión y de la
forma y orientación del objeto; y (iii) la reflexión especular de la propia
superficie, que en la configuración biestática llega al receptor y compite
con el eco del blanco.

\section{Escenario y geometría}
\label{sec:escenario}

Se considera una interfaz plana horizontal ($z=0$) entre el aire ($n_1=1$) y
un suelo homogéneo con índice de refracción $n_2=2$ (permitividad relativa
$\varepsilon_r=4$). El blanco se sitúa con su centro en
$\vect{p}_t=(0,0,-2)\,\mathrm{m}$.

El transmisor recorre un arco de radio constante $r=100\,\mathrm{m}$
centrado en el punto de la superficie situado sobre el blanco, contenido en
el plano $xz$:
\begin{equation}
  \vect{p}_{Tx}(\epsilon) = r\,(\cos\epsilon,\,0,\,\sin\epsilon),
  \qquad
  \vect{p}_{Rx}^{B}(\epsilon) = r\,(-\cos\epsilon,\,0,\,\sin\epsilon),
  \label{eq:geometria}
\end{equation}
donde $\epsilon$ es la elevación sobre el horizonte ($\epsilon=90^\circ$
corresponde al nadir). La elevación se discretiza en
$\epsilon=20^\circ,25^\circ,\dots,90^\circ$ (15 posiciones). En el nadir
ambas configuraciones coinciden, por lo que las comparaciones agregadas se
restringen a las 14 posiciones restantes.

Para cada par emisor--blanco se determina el punto de refracción en la
interfaz mediante el principio de Fermat, que para una interfaz plana se
reduce a resolver la ley de Snell,
\begin{equation}
  n_1\sin\theta_i = n_2\sin\theta_t ,
\end{equation}
donde $\theta_i$ es el ángulo de incidencia en el aire, medido desde la
normal, y $\theta_t$ el de transmisión en el suelo. Como el punto de
refracción se desplaza menos de 1\,m del origen, $\theta_i$ difiere de
$90^\circ-\epsilon$ en menos de $0.31^\circ$ en todo el barrido. La
elevación $\epsilon$ recorre, por tanto, incidencias entre
$\theta_i\approx 70^\circ$ y $0^\circ$, e incluye el ángulo de Brewster
\begin{equation}
  \theta_B = \arctan\frac{n_2}{n_1} = 63.43^\circ ,
\end{equation}
que en el arco corresponde a $\epsilon_B = 26.31^\circ$ y en el suelo a
$\theta_{t,B}=26.57^\circ$. Por la ley de Snell, desde el aire solo son
accesibles direcciones de observación dentro del cono
$\theta_t\le\theta_c=\arcsin(n_1/n_2)=30^\circ$ en torno a la vertical del
blanco.

La Tabla~\ref{tab:parametros} resume los parámetros del sistema, que
corresponden a la configuración de referencia de las simulaciones previas de
la tesis, con una frecuencia portadora de 425\,MHz.

\begin{table}[H]
\centering
\caption{Parámetros del sistema y del escenario.}
\label{tab:parametros}
\begin{tabular}{@{}llr@{}}
\toprule
\textbf{Magnitud} & \textbf{Símbolo} & \textbf{Valor} \\
\midrule
Frecuencia portadora & $f_c$ & 425\,MHz \\
Longitud de onda en el aire / en el suelo & $\lambda_0$ / $\lambda_s$ & 0.706\,m / 0.353\,m \\
Potencia transmitida & $P_t$ & 100\,kW \\
Ganancia de las antenas Tx y Rx & $G_t$, $G_r$ & 6\,dBi \\
Ancho de banda & $B$ & 50\,MHz \\
Temperatura de ruido & $T$ & 300\,K \\
Potencia de ruido térmico & $kTB$ & $-96.84\dBm$ \\
Índice de refracción del suelo & $n_2$ & 2 \\
Tangente de pérdidas del suelo (no incluida) & $\tan\delta$ & 0.05 \\
Radio del arco de vuelo & $r$ & 100\,m \\
Profundidad del centro del blanco & $d$ & 2\,m \\
Elevaciones simuladas & $\epsilon$ & $20^\circ$:$5^\circ$:$90^\circ$ \\
Polarización (Tx y Rx) & --- & TM (vertical) \\
\bottomrule
\end{tabular}
\end{table}

\section{Modelo de propagación en la interfaz}
\label{sec:fresnel}

La transmisión a través de la interfaz se modela con los coeficientes de
Fresnel para polarización TM (campo eléctrico en el plano de incidencia),
la única que presenta ángulo de Brewster:
\begin{equation}
  r_{TM} = \frac{n_2\cos\theta_i - n_1\cos\theta_t}{n_2\cos\theta_i + n_1\cos\theta_t},
  \qquad
  t_{TM} = \frac{2n_1\cos\theta_i}{n_2\cos\theta_i + n_1\cos\theta_t}.
\end{equation}
La transmitancia y la reflectancia de potencia son
\begin{equation}
  T = \frac{n_2\cos\theta_t}{n_1\cos\theta_i}\,|t_{TM}|^2,
  \qquad
  \Gamma = |r_{TM}|^2,
  \qquad
  T+\Gamma = 1 .
\end{equation}
El trayecto de ida (aire$\to$suelo) tiene transmitancia $T_\downarrow$ y el
de vuelta (suelo$\to$aire) $T_\uparrow$; por reciprocidad, para el mismo par
de ángulos $T_\uparrow=T_\downarrow$. En la configuración especular, la
simetría de la geometría hace que ambos trayectos compartan también el mismo
ángulo, de modo que el producto $T_\downarrow T_\uparrow$ es idéntico para
M y B en cada posición (verificado numéricamente con un error máximo de
$4\times10^{-16}$).

\section{Ecuación radar}
\label{sec:radar}

La potencia recibida del blanco se calcula con la ecuación radar biestática,
incluyendo la transmitancia de ambos trayectos:
\begin{equation}
  P_r = \frac{P_t\,G_t\,G_r\,\lambda_0^2\,\sigma(\uv{\imath},\uv{s})\,
              T_\downarrow T_\uparrow}
             {(4\pi)^3\,R_T^2\,R_R^2},
  \label{eq:radar}
\end{equation}
donde $R_T$ y $R_R$ son las longitudes de los trayectos Tx--blanco y
blanco--Rx (tramo en aire más tramo en suelo), y $\sigma$ es la RCS del
blanco evaluada en el medio (Sección~\ref{sec:rcs}), función de la
dirección hacia el transmisor $\uv{\imath}$ y hacia el receptor $\uv{s}$,
ambas definidas dentro del suelo a partir de los rayos refractados. En el
caso monoestático $\uv{s}=\uv{\imath}$ y $R_R=R_T$. Las antenas se suponen
apuntadas al punto de refracción, de modo que operan con su ganancia
máxima. El ruido térmico de referencia es $P_n=kTB$, sin figura de ruido ni
integración coherente.

En la configuración biestática especular, el receptor recibe además la
reflexión especular de la superficie. Por teoría de imágenes, para una
interfaz plana e infinita,
\begin{equation}
  P_{sup} = \frac{P_t\,G_t\,G_r\,\lambda_0^2\,\Gamma(\theta_i)}
                 {(4\pi)^2\,\left(|\vect{p}_{Tx}|+|\vect{p}_{Rx}|\right)^2}.
  \label{eq:superficie}
\end{equation}
Como $\Gamma(\theta_B)=0$, esta contribución se anula en el ángulo de
Brewster. Se define la relación blanco--superficie
$\mathrm{SCR}=P_r^B/P_{sup}$ para cuantificar esta competencia.

\section{Modelo de sección eficaz radar biestática}
\label{sec:rcs}

\subsection{Óptica Física biestática}

La RCS de los blancos, conductores eléctricos perfectos (PEC), se calcula
con la aproximación de Óptica Física (PO). La corriente superficial inducida
en las caras iluminadas por el transmisor es
$\vect{J}_s=2\,\uv{n}\times\vect{H}_i$, y el campo dispersado en campo lejano
en la dirección $\uv{s}$, proyectado sobre la polarización de recepción
$\uv{e}_r$, conduce a
\begin{equation}
  \sigma(\uv{\imath},\uv{s}) = \frac{4\pi}{\lambda_s^2}
  \left|\,\sum_{m:\,\uv{n}_m\cdot\uv{\imath}>0}
  \uv{e}_r\cdot\big[\uv{n}_m\times(\uv{k}_i\times\uv{e}_i)\big]
  \int_{S_m} e^{\,j k_s(\uv{\imath}+\uv{s})\cdot\vect{r}'}\,dS'\right|^2 ,
  \label{eq:po}
\end{equation}
donde $\uv{k}_i=-\uv{\imath}$ es la dirección de propagación incidente,
$\uv{e}_i$ y $\uv{e}_r$ son los vectores de polarización TM de transmisión
y recepción (contenidos en el plano de incidencia $xz$), $\uv{n}_m$ es la
normal saliente de la cara $m$ y $k_s=2\pi/\lambda_s$. La superficie de cada
objeto se representa con una malla de triángulos planos, sobre los que la
integral de fase se evalúa en forma cerrada y exacta
\cite{gordon1975}; así, las caras planas no necesitan subdividirse. En el
caso monoestático ($\uv{s}=\uv{\imath}$), el factor de polarización se
reduce a $\uv{n}_m\cdot\uv{\imath}$ y la expresión coincide con la
formulación PO monoestática clásica \cite{knott1993}.

Como el blanco está inmerso en el suelo, la RCS se evalúa con la longitud de
onda del medio, $\lambda_s=\lambda_0/n_2=0.353\,\mathrm{m}$ (equivalente a
una frecuencia en el vacío de 850\,MHz), y en las direcciones de los rayos
refractados ($\theta_t$), no en las del aire.

\subsection{Objetos simulados}

Se simulan los cinco sólidos cerrados de la biblioteca de RCS del grupo,
con geometría centrada en el origen del blanco y ejes de revolución
verticales (Tabla~\ref{tab:objetos}). Las superficies curvas se facetan con
un número de segmentos por vuelta $N$ tal que la flecha de cada faceta es
menor que $\lambda_s/100$, $a\,[1-\cos(\pi/N)]\le\lambda_s/100$, criterio
que acota el error de fase de la faceta. Todos los objetos tienen una
dimensión máxima de 1\,m ($2.83\,\lambda_s$). La geometría simulada de cada
uno se muestra en el Apéndice~\ref{ap:figuras}
(Figuras~\ref{fig:cubo_geometria} y \ref{fig:caja_geometria}--\ref{fig:cono_geometria}).

\begin{table}[H]
\centering
\caption{Objetos PEC simulados. $\sigma_{nadir}$: RCS monoestática en
incidencia vertical calculada con PO; referencia: valor de óptica
geométrica/física canónico.}
\label{tab:objetos}
\small
\setlength{\tabcolsep}{4pt}
\begin{tabular}{@{}llrrr@{}}
\toprule
\textbf{Objeto} & \textbf{Dimensiones} & \textbf{Triángulos} & $\bm{\sigma_{nadir}}$ \textbf{(PO)} & \textbf{Referencia} \\
\midrule
Cubo & arista 1\,m & 12 & $20.04\dBsm$ & $4\pi A^2/\lambda_s^2 = 20.04\dBsm$ \\
Caja & $1\times0.5\times0.3$\,m & 12 & $14.02\dBsm$ & $4\pi A^2/\lambda_s^2 = 14.02\dBsm$ \\
Cilindro & $a=0.3$\,m, $L=1$\,m, eje $z$ & 84 & $8.94\dBsm$ & tapa: $9.06\dBsm$ \\
Esfera & $a=0.5$\,m & 702 & $-1.54\dBsm$ & $\pi a^2 = -1.05\dBsm$ \\
Cono & base $a=0.3$\,m, $L=1$\,m, vértice arriba & 42 & $-15.86\dBsm$ & --- (vértice) \\
\bottomrule
\end{tabular}
\end{table}

\subsection{Validación del modelo de RCS}
\label{sec:validacion_met}

El modelo se validó por tres vías independientes:
\begin{enumerate}[itemsep=2pt]
  \item Comparación del caso monoestático con resultados previos de la
  biblioteca de RCS del grupo, calculados con PO y con el Método de los
  Momentos (MoM, ecuación integral de campo eléctrico) a 400\,MHz, en el
  corte de elevación del plano de incidencia.
  \item Comparación del caso biestático especular de la cara superior del
  cubo con la expresión analítica de la placa,
  $\sigma=4\pi A^2\cos^2\theta/\lambda^2$.
  \item Verificación de que evaluar un objeto rotado con geometría fija
  equivale a evaluar el objeto sin rotar con todas las direcciones y
  polarizaciones rotadas por $\mathbf{R}^{\mathsf T}$ (invariancia de
  \eqref{eq:po} bajo rotaciones rígidas), propiedad que se usa en el
  Monte Carlo.
\end{enumerate}

\section{Diagramas de dispersión tridimensionales}

Para cada objeto se calcula la RCS sobre la esfera completa de direcciones
de observación (paso de $0.5^\circ$ en acimut y elevación) en dos casos:
monoestático ($\uv{s}=\uv{\imath}$ recorriendo la esfera) y biestático con
el transmisor fijo en la dirección refractada de Brewster
($\theta_{t,B}=26.57^\circ$, lado $+x$) y el receptor recorriendo la esfera.
Ambos diagramas se representan con un rango dinámico de 30\,dB. Además, en
cortes de $0.05^\circ$ se mide el semiancho a $-3\dB$ del lóbulo especular
biestático, que determina la tolerancia de posicionamiento del receptor.

\section{Experimento de Monte Carlo sobre la orientación del blanco}
\label{sec:mc}

\subsection{Diseño}

La orientación de un blanco enterrado es, en general, desconocida. Para
evaluar la robustez de la comparación M/B frente a esta incertidumbre se
realiza un experimento de Monte Carlo: para \emph{cada} una de las 15
posiciones del barrido se generan $N=10^4$ orientaciones aleatorias e
independientes del objeto, rotándolo en torno a su centro con
\begin{equation}
  \mathbf{R} = \mathbf{R}_y(\beta_y)\,\mathbf{R}_x(\beta_x),
  \qquad \beta_x,\beta_y \sim \mathcal{U}(-180^\circ,180^\circ),
\end{equation}
es decir, giros aleatorios independientes alrededor de los dos ejes
horizontales, con un rango que permite a cualquier objeto adoptar cualquier
inclinación. En cada repetición se calcula la potencia recibida en M y en B
\emph{con el mismo objeto}, lo que convierte la comparación en un diseño
pareado. La geometría, las transmitancias y los factores de propagación de
\eqref{eq:radar} son los de la orientación nominal; solo cambia $\sigma$.
Por la invariancia bajo rotaciones, cada evaluación se realiza con el
objeto sin rotar y los vectores
$\mathbf{R}^{\mathsf T}\uv{\imath}$, $\mathbf{R}^{\mathsf T}\uv{s}$,
$\mathbf{R}^{\mathsf T}\uv{e}_i$ y $\mathbf{R}^{\mathsf T}\uv{e}_r$, con un
error numérico inferior a $4\times10^{-14}\dB$ respecto de rotar la malla.

\subsection{Estimadores}

Sea $P_{r,k}^{M}(n)$ la potencia recibida en la posición $k$ y la
repetición $n$. Se emplean los siguientes estimadores:
\begin{itemize}[itemsep=3pt]
  \item \textbf{Potencia media (energía):}
  $\bar P_k^{M}=\frac1N\sum_n P_{r,k}^{M}(n)$, en unidades lineales; es la
  magnitud que determina la energía útil acumulada en la integración SAR.
  \item \textbf{Diferencia de medias:}
  $\Delta_k = 10\log_{10}\left(\bar P_k^{B}/\bar P_k^{M}\right)$; y su
  versión agregada sobre las 14 posiciones fuera del nadir,
  $\Delta = 10\log_{10}\left(\sum_{k}\sum_n P_{r,k}^{B}(n) \big/
  \sum_{k}\sum_n P_{r,k}^{M}(n)\right)$.
  \item \textbf{Intervalos de confianza:} intervalo del 95\,\% por
  \emph{bootstrap} de percentiles pareado \cite{efron1993}, con 2000
  remuestreos de las repeticiones (las mismas repeticiones para M y B).
  Una posición se declara ``significativamente mejor/peor'' si el intervalo
  de $\Delta_k$ excluye el cero.
  \item \textbf{Mediana y probabilidad de superioridad:} la mediana de
  $P_{r,k}$ describe el caso típico, y $\Pr\{P_{r,k}^{B}>P_{r,k}^{M}\}$ la
  fracción de repeticiones en que B recibe más que M.
  \item \textbf{Aciertos especulares:} número de repeticiones en que
  $\sigma$ supera el máximo monoestático del objeto menos 10\,dB; indica
  cuántas veces una cara queda orientada en reflexión especular, eventos
  que dominan la media y de los que depende su precisión estadística.
\end{itemize}

\section{Hipótesis y limitaciones del modelo}
\label{sec:limitaciones}

\begin{enumerate}[itemsep=2pt]
  \item Blancos PEC; suelo homogéneo, isótropo y sin pérdidas (con
  $\tan\delta=0.05$, las pérdidas dieléctricas añadirían unos
  $3.9\dB/\mathrm{m}$ de trayecto en el suelo, del orden de $15\dB$ en el
  nadir, idénticos para M y B).
  \item Interfaz plana y lisa; no se modela la rugosidad ni el
  \emph{clutter} volumétrico.
  \item PO no incluye difracción en aristas y vértices, ondas de
  superficie ni reflexiones múltiples. Con objetos de $2.83\,\lambda_s$ y
  detalles menores que $\lambda_s$ (tapa del cilindro, vértice del cono),
  los valores absolutos fuera de las direcciones especulares tienen una
  incertidumbre de varios dB.
  \item Polarización TM co-polar; no se calculan TE ni la polarización
  cruzada.
  \item Aproximación de campo lejano para el blanco y ecuación radar con
  distancia geométrica (tramo en aire más tramo en suelo), sin el factor de
  divergencia por refracción; esta simplificación afecta por igual a M y B.
  \item Análisis de onda continua a una sola frecuencia, sin promediado en
  la banda de la señal.
\end{enumerate}

% ══════════════════════════════════════════════════════════════════════════
\chapter{Resultados}
\label{cap:resultados}

\section{Validación del modelo de RCS}
\label{sec:validacion}

La Tabla~\ref{tab:validacion} resume la validación. La formulación PO
reproduce los resultados PO de referencia con una diferencia media de
$0.01\dB$, lo que confirma la implementación de la integral exacta por
triángulo y la normalización. Frente a MoM, que incluye todos los
mecanismos de onda completa, las diferencias medias son de $2.3\dB$ para el
cubo y $4.4\dB$ para el cilindro, en los lóbulos situados hasta 20\,dB por
debajo del máximo. Estas diferencias son las esperables para objetos de
1.3 longitudes de onda, en el límite inferior de validez de PO; en el
caso del cilindro, la tapa ($a=0.4\,\lambda$) está en la región resonante y
PO subestima su retorno vertical en 3\,dB. La expresión biestática
reproduce la placa analítica con diferencias de $0.01$--$0.09\dB$ hasta
$45^\circ$; las pequeñas desviaciones se deben a la contribución de las
caras laterales del cubo, ausente en el modelo de placa aislada.

\begin{table}[H]
\centering
\caption{Validación del modelo PO. Frecuencias de 400\,MHz (en el vacío)
para las referencias de la biblioteca de RCS y de 850\,MHz equivalentes
($\lambda_s$) para la placa biestática.}
\label{tab:validacion}
\small
\setlength{\tabcolsep}{4pt}
\begin{tabular}{@{}lllrrr@{}}
\toprule
\textbf{Caso} & \textbf{Referencia} & \textbf{Tamaño} & \textbf{PO} & \textbf{Ref.} & \textbf{Error medio} \\
\midrule
Cubo 3\,m, monoestático & PO & $4.0\,\lambda$ & $32.58\dBsm$ & $32.58\dBsm$ & $0.01\dB$ \\
Cubo 1\,m, monoestático & MoM & $1.3\,\lambda$ & $13.49\dBsm$ & $12.67\dBsm$ & $2.27\dB$ \\
Cilindro $0.3\times1$\,m, monoestático & MoM & $1.3\,\lambda$ & $2.26\dBsm$ & $5.31\dBsm$ & $4.38\dB$ \\
\midrule
Placa 1\,m, biestático $\theta=0^\circ$ & Analítica & $2.8\,\lambda_s$ & $20.038\dBsm$ & $20.038\dBsm$ & $0.000\dB$ \\
Placa 1\,m, biestático $\theta=15^\circ$ & Analítica & & $19.727\dBsm$ & $19.737\dBsm$ & $0.010\dB$ \\
Placa 1\,m, biestático $\theta=30^\circ$ & Analítica & & $18.878\dBsm$ & $18.789\dBsm$ & $0.089\dB$ \\
Placa 1\,m, biestático $\theta=45^\circ$ & Analítica & & $17.013\dBsm$ & $17.028\dBsm$ & $0.015\dB$ \\
\bottomrule
\end{tabular}
\end{table}

\section{Efecto de la interfaz: transmitancia y reflexión de superficie}
\label{sec:res_fresnel}

La transmitancia de ida y vuelta $T_\downarrow T_\uparrow$ alcanza la unidad
en $\epsilon=26.30^\circ$ ($\theta_i=63.44^\circ$), frente a $0.790$ en el
nadir ($-1.02\dB$). La \textbf{ganancia de transmitancia atribuible al
ángulo de Brewster es, por tanto, de $+1.02\dB$} respecto de la incidencia
vertical. Para un blanco puntual de RCS constante ($\sigma=1\,\mathrm{m}^2$),
la potencia recibida en $\epsilon_B$ supera a la del nadir en $+1.14\dB$
($+1.02\dB$ de transmitancia y $+0.11\dB$ por la ligera reducción del
trayecto), idéntica para M y B. La posición del barrido más cercana a
Brewster ($\epsilon=25^\circ$, $\theta_i=64.75^\circ$) conserva
$T_\downarrow T_\uparrow=0.999$.

La reflectancia de superficie $\Gamma$ varía entre $0.111$ en el nadir y
cero en Brewster. En consecuencia, la reflexión especular de la superficie
que llega al receptor biestático varía entre $+11.4\dBm$ en el nadir y
$-12.1\dBm$ en $\epsilon=25^\circ$, anulándose en $\epsilon_B$
(Figura~\ref{fig:cubo_principal}, panel inferior izquierdo). Esta es la
contribución que más depende del ángulo en todo el balance: en el barrido
continuo varía unos $67\dB$, de los cuales $44\dB$ se concentran en los
$1.3^\circ$ que separan $\epsilon=25^\circ$ de $\epsilon_B$.

\section{Orientación nominal}
\label{sec:res_nominal}

\subsection{Cubo: análisis detallado}

La Figura~\ref{fig:cubo_geometria} muestra la geometría simulada: un cubo
PEC de 1\,m de arista, descrito exactamente por 12 triángulos, con su centro
a 2\,m de profundidad y las caras alineadas con los ejes.

\begin{figure}[H]
\centering
\includegraphics[width=0.92\textwidth]{fig_cubo_geometria.png}
\caption{Geometría simulada del cubo de 1\,m. (a) Malla de facetas en
coordenadas del objeto. (b) Sección en el plano de incidencia $xz$ con el
objeto enterrado; la línea discontinua indica la profundidad del centro.}
\label{fig:cubo_geometria}
\end{figure}

La Tabla~\ref{tab:cubo_nominal} y la Figura~\ref{fig:cubo_principal}
presentan el balance del cubo en orientación nominal (caras alineadas con
los ejes). La RCS biestática especular permanece casi constante en todo el
barrido ($18.95$--$20.04\dBsm$), mientras que la monoestática varía entre
$-16.1$ y $+20.0\dBsm$ y solo alcanza el máximo en el nadir. La potencia
biestática es, por tanto, prácticamente independiente de la elevación
($-5.2$ a $-5.4\dBm$), y la ganancia biestática respecto de la
monoestática oscila entre $+0.9$ y $+35.5\dB$ (media aritmética de
$+16.6\dB$ fuera del nadir). En términos energéticos, sumando las 14
posiciones fuera del nadir, la configuración especular recibe
\textbf{$+9.56\dB$ más energía} que la monoestática.

\begin{table}[H]
\centering
\caption{Cubo de 1\,m en orientación nominal: balance por posición.
$\beta$: ángulo biestático en el suelo. SCR: relación entre la potencia
biestática del blanco y la reflexión especular de la superficie.}
\label{tab:cubo_nominal}
\small
\begin{tabular}{@{}rrrrrrrrrrr@{}}
\toprule
$\epsilon$ & $\theta_i$ & $\theta_t$ & $\beta$ & $T_\downarrow T_\uparrow$ & $\sigma^M$ & $\sigma^B$ & $P_r^M$ & $P_r^B$ & $\Delta$ & SCR \\
{[°]} & {[°]} & {[°]} & {[°]} & [dB] & [dBsm] & [dBsm] & [dBm] & [dBm] & [dB] & [dB] \\
\midrule
20 & 69.79 & 27.98 & 55.97 & $-0.13$ & $-0.63$ & 18.95 & $-24.98$ & $-5.40$ & 19.58 & $-8.10$ \\
25 & 64.75 & 26.89 & 53.77 & $-0.00$ & 0.76 & 19.00 & $-23.48$ & $-5.24$ & 18.24 & 6.85 \\
30 & 59.72 & 25.58 & 51.16 & $-0.03$ & 1.66 & 19.06 & $-22.61$ & $-5.21$ & 17.40 & $-1.10$ \\
35 & 54.70 & 24.08 & 48.17 & $-0.12$ & 1.36 & 19.14 & $-23.02$ & $-5.23$ & 17.78 & $-7.59$ \\
40 & 49.69 & 22.41 & 44.83 & $-0.24$ & $-1.76$ & 19.25 & $-26.27$ & $-5.26$ & 21.01 & $-10.66$ \\
45 & 44.69 & 20.59 & 41.18 & $-0.38$ & $-16.09$ & 19.37 & $-40.75$ & $-5.28$ & 35.47 & $-12.53$ \\
50 & 39.70 & 18.63 & 37.25 & $-0.51$ & $-0.10$ & 19.51 & $-24.90$ & $-5.29$ & 19.61 & $-13.78$ \\
55 & 34.72 & 16.55 & 33.09 & $-0.62$ & 5.66 & 19.64 & $-19.27$ & $-5.29$ & 13.98 & $-14.67$ \\
60 & 29.75 & 14.36 & 28.73 & $-0.73$ & 7.03 & 19.76 & $-18.01$ & $-5.28$ & 12.73 & $-15.31$ \\
65 & 24.78 & 12.10 & 24.19 & $-0.82$ & 3.63 & 19.86 & $-21.51$ & $-5.28$ & 16.23 & $-15.80$ \\
70 & 19.81 & 9.76 & 19.52 & $-0.89$ & $-6.92$ & 19.93 & $-32.15$ & $-5.29$ & 26.86 & $-16.16$ \\
75 & 14.86 & 7.37 & 14.73 & $-0.95$ & 10.55 & 19.99 & $-14.74$ & $-5.30$ & 9.44 & $-16.42$ \\
80 & 9.90 & 4.93 & 9.87 & $-0.99$ & 16.35 & 20.02 & $-8.98$ & $-5.31$ & 3.67 & $-16.61$ \\
85 & 4.95 & 2.47 & 4.95 & $-1.02$ & 19.16 & 20.04 & $-6.20$ & $-5.32$ & 0.88 & $-16.72$ \\
90 & 0.00 & 0.00 & 0.00 & $-1.02$ & 20.04 & 20.04 & $-5.33$ & $-5.33$ & 0.00 & $-16.76$ \\
\bottomrule
\end{tabular}
\end{table}

\begin{figure}[H]
\centering
\includegraphics[width=\textwidth]{fig_cubo_principal.png}
\caption{Cubo de 1\,m en orientación nominal. Fila superior: geometría del
arco en el plano $xz$ (el trazo discontinuo marca la posición más cercana a
Brewster), coeficientes de Fresnel TM y RCS monoestática y biestática en
función de la elevación. Fila inferior: potencia recibida (incluida la
reflexión especular de la superficie), ganancia biestática por posición y
relación blanco--superficie en el receptor biestático. Las curvas continuas
corresponden a un barrido de $0.05^\circ$ y los marcadores a las posiciones
simuladas.}
\label{fig:cubo_principal}
\end{figure}

La relación blanco--superficie (SCR) del receptor biestático es negativa en
casi todo el barrido (entre $-8$ y $-17\dB$): con la superficie lisa, la
reflexión de la interfaz supera al eco del cubo. Solo en el entorno de
Brewster se invierte, alcanzando $+6.9\dB$ en $\epsilon=25^\circ$ y un
máximo de $50.6\dB$ en $\epsilon_B$ en el barrido continuo. Este resultado
muestra que, en la configuración especular, el papel decisivo del ángulo de
Brewster no es aumentar la energía que llega al blanco ($+1\dB$), sino
suprimir la reflexión de superficie que la enmascara. El intervalo útil es
estrecho: a $1.3^\circ$ de $\epsilon_B$ la SCR ya ha caído unos $44\dB$.

\subsection{Todos los objetos}

La Tabla~\ref{tab:nominal} resume la orientación nominal para los cinco
objetos. Los tres objetos con una cara plana horizontal en la parte
superior ---cubo, caja y cilindro, cuya tapa superior es un disco--- muestran
el mismo mecanismo: la RCS biestática especular se mantiene próxima al
retorno vertical de esa cara en todo el barrido, mientras que la
monoestática cae decenas de dB fuera del nadir. La ganancia energética del
modo especular es de $+9.6\dB$ (cubo y caja) y $+6.2\dB$ (cilindro). La
esfera, sin caras planas, tiene una RCS casi isótropa en ambos modos y la
diferencia es pequeña ($-0.5\dB$). El cono, con el vértice hacia arriba, no
presenta cara especular accesible: su RCS observable en ambos modos es
débil ($-31$ a $-16\dBsm$), asociada a las contribuciones del vértice y del
borde de la base, que PO solo representa de forma aproximada; la diferencia
neta es de $+2.0\dB$, con valores por posición de signo variable.

\begin{table}[H]
\centering
\caption{Orientación nominal: rangos de RCS en el barrido y diferencia de
potencia biestática especular menos monoestática. ``Energía'': diferencia
de la potencia sumada en las 14 posiciones fuera del nadir. ``$\epsilon=25^\circ$'':
posición más cercana a Brewster.}
\label{tab:nominal}
\begin{tabular}{@{}lrrrrrr@{}}
\toprule
& \multicolumn{2}{c}{\textbf{RCS [dBsm]}} & \multicolumn{3}{c}{$\bm{\Delta = P_r^B - P_r^M}$ \textbf{por posición [dB]}} & \textbf{Energía} \\
\cmidrule(lr){2-3}\cmidrule(lr){4-6}
\textbf{Objeto} & \textbf{Mono} & \textbf{Biest.} & \textbf{mín.} & \textbf{máx.} & $\bm{\epsilon=25^\circ}$ & \textbf{[dB]} \\
\midrule
Cubo     & $-16.1$ a $20.0$ & $19.0$ a $20.0$ & $+0.9$ & $+35.5$ & $+18.2$ & $+9.56$ \\
Caja     & $-19.3$ a $14.0$ & $12.9$ a $14.0$ & $+0.9$ & $+32.8$ & $+16.3$ & $+9.56$ \\
Cilindro & $-23.4$ a $8.9$  & $7.8$ a $9.0$   & $+0.3$ & $+31.5$ & $+17.7$ & $+6.17$ \\
Esfera   & $-1.5$ a $-0.7$  & $-3.0$ a $-0.1$ & $-2.0$ & $+0.9$  & $+0.2$  & $-0.54$ \\
Cono     & $-23.7$ a $-15.9$ & $-31.4$ a $-15.9$ & $-10.0$ & $+7.7$ & $-9.0$ & $+2.00$ \\
\bottomrule
\end{tabular}
\end{table}

\subsection{Desglose por posición}
\label{sec:res_nominal_angulo}

La cifra energética de la Tabla~\ref{tab:nominal} agrega las 14 posiciones
y oculta en qué parte del barrido es mejor cada configuración. La
Tabla~\ref{tab:nominal_posicion} y la Figura~\ref{fig:resumen_angulo}(a)
desglosan la diferencia $\Delta_k$ por objeto y por posición; los balances
completos de cada objeto (RCS, potencias y SCR) están en el
Apéndice~\ref{ap:figuras}. Se distinguen tres comportamientos:
\begin{itemize}[itemsep=2pt]
  \item \textbf{Cubo, caja y cilindro:} el especular es superior en
  \emph{todas} las posiciones. La ventaja es de $+16$ a $+20\dB$ en las
  incidencias rasantes y en el entorno de Brewster ($20^\circ$--$35^\circ$) y
  decrece hacia el nadir, donde ambos modos coinciden. Los picos de
  $\epsilon=45^\circ$ ($+31.5$ a $+35.5\dB$) y $70^\circ$ ($+27\dB$ en cubo y
  caja) no son ganancias del modo especular, sino nulos de la RCS
  monoestática entre lóbulos. En el cilindro la ventaja cae antes
  ($<9\dB$ a partir de $60^\circ$) porque el lóbulo de la tapa, más ancho,
  entra antes en el campo de visión monoestático.
  \item \textbf{Esfera:} diferencias pequeñas de signo variable; el especular
  es ligeramente mejor entre $25^\circ$ y $40^\circ$ (hasta $+0.9\dB$) y el
  monoestático entre $45^\circ$ y $85^\circ$ (hasta $-2.0\dB$).
  \item \textbf{Cono:} el signo se invierte a lo largo del barrido. El
  monoestático es mejor en las incidencias rasantes, incluida la posición de
  Brewster ($-10.0$ a $-1.0\dB$ entre $20^\circ$ y $45^\circ$), y el
  especular a partir de $50^\circ$ (hasta $+7.7\dB$ en $65^\circ$). El
  balance energético positivo ($+2.0\dB$) se debe a estas últimas
  posiciones.
\end{itemize}

\begin{table}[H]
\centering
\caption{Orientación nominal: diferencia de potencia $\Delta_k=P_r^B-P_r^M$
[dB] por posición y objeto. Valores positivos: el especular recibe más. En
el nadir ($\epsilon=90^\circ$) la diferencia es nula por construcción.}
\label{tab:nominal_posicion}
\begin{tabular}{@{}rrrrrrr@{}}
\toprule
$\epsilon$ [°] & $\theta_i$ [°] & \textbf{Cubo} & \textbf{Caja} & \textbf{Cilindro} & \textbf{Esfera} & \textbf{Cono} \\
\midrule
20 & 69.79 & $+19.58$ & $+17.35$ & $+17.66$ & $-0.15$ & $-10.01$ \\
25 & 64.75 & $+18.24$ & $+16.30$ & $+17.67$ & $+0.17$ & $-9.05$ \\
30 & 59.72 & $+17.40$ & $+16.11$ & $+18.05$ & $+0.89$ & $-8.12$ \\
35 & 54.70 & $+17.78$ & $+17.40$ & $+19.52$ & $+0.26$ & $-2.41$ \\
40 & 49.69 & $+21.01$ & $+21.83$ & $+23.79$ & $+0.49$ & $-1.89$ \\
45 & 44.69 & $+35.47$ & $+32.82$ & $+31.52$ & $-0.37$ & $-1.00$ \\
50 & 39.70 & $+19.61$ & $+18.87$ & $+18.82$ & $-2.03$ & $+2.80$ \\
55 & 34.72 & $+13.98$ & $+14.00$ & $+12.51$ & $-1.97$ & $+5.81$ \\
60 & 29.75 & $+12.73$ & $+12.95$ & $+8.53$ & $-1.75$ & $+7.43$ \\
65 & 24.78 & $+16.23$ & $+16.52$ & $+5.70$ & $-1.65$ & $+7.65$ \\
70 & 19.81 & $+26.86$ & $+26.64$ & $+3.60$ & $-1.69$ & $+6.24$ \\
75 & 14.86 & $+9.44$ & $+9.45$ & $+2.04$ & $-1.27$ & $+3.84$ \\
80 & 9.90 & $+3.67$ & $+3.68$ & $+0.93$ & $-0.58$ & $+1.75$ \\
85 & 4.95 & $+0.88$ & $+0.88$ & $+0.25$ & $-0.18$ & $+0.44$ \\
\midrule
\multicolumn{2}{@{}l}{Energía (14 pos.)} & $+9.56$ & $+9.56$ & $+6.17$ & $-0.54$ & $+2.00$ \\
\bottomrule
\end{tabular}
\end{table}

\section{Diagramas de dispersión tridimensionales}
\label{sec:res_3d}

Las Figuras~\ref{fig:rcs3d_cubo}--\ref{fig:rcs3d_cono} muestran los
diagramas 3D de RCS de cada objeto en el suelo, y la
Tabla~\ref{tab:rcs3d} sus valores característicos. El diagrama
monoestático del cubo presenta seis lóbulos especulares, uno por cara, de
los cuales solo el vertical superior cae dentro del cono de direcciones
accesibles desde el aire ($\theta_t\le30^\circ$). Con el transmisor fijo en
la dirección de Brewster, el diagrama biestático concentra la energía en
dos lóbulos: el especular de la cara superior, orientado hacia el receptor
simétrico ($19.0\dBsm$), y el de dispersión hacia delante, en la sombra del
objeto, que es el máximo absoluto ($22.6\dBsm$) pero apunta hacia abajo y
no es observable desde el aire. En la misma dirección de Brewster, el
retorno monoestático del cubo es de solo $1.1\dBsm$: la diferencia de
$18\dB$ explica la ganancia del modo especular en orientación nominal.

\begin{table}[H]
\centering
\caption{Valores característicos de los diagramas de RCS con el transmisor
en la dirección refractada de Brewster ($\theta_{t,B}=26.57^\circ$).
Semiancho: semiancho a $-3\dB$ del lóbulo especular biestático en el suelo.
Máx. biest.: máximo del diagrama biestático, en la dirección de
dispersión hacia delante. Tolerancia: desplazamiento admisible del receptor
sobre el arco de 100\,m, hacia el nadir, con pérdida inferior a 3\,dB.}
\label{tab:rcs3d}
\small
\begin{tabular}{@{}lrrrrr@{}}
\toprule
\textbf{Objeto} & $\bm{\sigma^M(\theta_{t,B})}$ & $\bm{\sigma^B}$ \textbf{especular} & \textbf{Máx. biest.} & \textbf{Semiancho} & \textbf{Tolerancia} \\
& [dBsm] & [dBsm] & [dBsm] & [°] & [m] \\
\midrule
Cubo     & 1.07   & 19.04  & 22.60 & 9.9  & 49 \\
Caja     & $-3.07$ & 13.09 & 14.27 & 9.9  & 50 \\
Cilindro & $-9.87$ & 7.90  & 14.42 & 19.9 & 87 \\
Esfera   & $-1.06$ & $-0.29$ & 17.84 & --- & --- \\
Cono     & $-21.35$ & $-28.57$ & 9.17 & --- & --- \\
\bottomrule
\end{tabular}
\end{table}

A 425\,MHz, el lóbulo especular es ancho (semiancho de $\approx10^\circ$ en
el suelo para las caras de 1\,m y de $\approx20^\circ$ para la tapa del
cilindro), de modo que el lóbulo se extiende hasta el ángulo crítico: hacia
incidencias rasantes cubre todas las direcciones accesibles, y hacia el
nadir el receptor puede desplazarse unos 50\,m (cubo, caja) u 87\,m
(cilindro) sobre el arco sin perder más de 3\,dB. La geometría especular es,
por tanto, robusta a errores de posicionamiento del receptor en esta banda.
En la esfera y el cono no existe un lóbulo especular definido en la
dirección simétrica, por lo que estas magnitudes no son aplicables. El
diagrama monoestático del cono muestra un lóbulo toroidal horizontal
(generatriz) y un lóbulo inferior (base), ambos fuera del cono accesible,
lo que explica su baja RCS observable.

\section{Monte Carlo: orientación aleatoria del blanco}
\label{sec:res_mc}

\subsection{Cubo}

La Figura~\ref{fig:cubo_mc} presenta el resultado del Monte Carlo para el
cubo. Con orientación aleatoria, la potencia media recibida por ambas
configuraciones es prácticamente la misma en todas las posiciones
($-16.3$ a $-17.7\dBm$), unos $11$--$12\dB$ por debajo del valor biestático
nominal. La diferencia de medias por posición $\Delta_k$ está comprendida
entre $-0.74$ y $+0.33\dB$; en 9 de las 14 posiciones su intervalo de
confianza contiene el cero, en 4 (las de menor elevación, $20^\circ$--$35^\circ$)
el modo monoestático es significativamente superior por menos de $0.8\dB$,
y en una ($75^\circ$) lo es el especular, por $0.3\dB$. La diferencia
agregada es $\Delta=-0.26\dB$ (IC\,95\,\%: $[-0.37,-0.16]\dB$) y la
probabilidad de que el especular reciba más en una repetición es $0.52$.
Entre 933 y 1104 repeticiones por posición (9--11\,\%) contienen un acierto
especular, suficiente para que la media esté bien determinada.

\begin{figure}[H]
\centering
\includegraphics[width=\textwidth]{fig_cubo_montecarlo.png}
\caption{Monte Carlo del cubo con $10^4$ orientaciones aleatorias por
posición. (a) Potencia media recibida con intervalo de confianza del 95\,\%
en ambas configuraciones; la línea punteada es la referencia de un blanco
puntual de $1\,\mathrm{m}^2$. (b) Diferencia de potencia media (especular
menos monoestática) con su intervalo de confianza pareado.}
\label{fig:cubo_mc}
\end{figure}

\subsection{Todos los objetos}

La Tabla~\ref{tab:mc_resumen} y la Figura~\ref{fig:resumen} resumen el Monte
Carlo para los cinco objetos, y la Tabla~\ref{tab:mc_posicion} detalla la
diferencia de medias por posición, que se analiza en la
Sección~\ref{sec:res_mc_angulo}. En todos los casos la diferencia
agregada es inferior a $0.4\dB$ en valor absoluto, frente a los $+6$ a
$+10\dB$ de la orientación nominal. Aunque, por el elevado número de
repeticiones, algunos intervalos de confianza excluyen el cero, las
magnitudes son despreciables frente a cualquier otra incertidumbre del
balance (en particular, frente al error de PO).

\begin{table}[H]
\centering
\caption{Monte Carlo ($10^4$ orientaciones aleatorias por posición, giros
en $x$ e $y$ uniformes en $\pm180^\circ$): diferencia energética agregada
especular menos monoestática en las 14 posiciones fuera del nadir,
probabilidad media de superioridad del especular, número de posiciones con
diferencia significativa y aciertos especulares mínimos por posición.}
\label{tab:mc_resumen}
\begin{tabular}{@{}lrrcc@{}}
\toprule
\textbf{Objeto} & $\bm{\Delta}$ \textbf{(IC 95\,\%) [dB]} & $\bm{\Pr\{B>M\}}$ & \textbf{Mejor / peor / igual} & \textbf{Aciertos} \\
\midrule
Cubo     & $-0.26$ \ $[-0.37,\,-0.16]$ & 0.52 & 1 / 4 / 9  & 933 \\
Caja     & $-0.18$ \ $[-0.26,\,-0.10]$ & 0.50 & 1 / 3 / 10 & 214 \\
Cilindro & $-0.26$ \ $[-0.30,\,-0.22]$ & 0.50 & 4 / 7 / 3  & 2330 \\
Esfera   & $-0.37$ \ $[-0.37,\,-0.36]$ & 0.38 & 4 / 10 / 0 & --- \\
Cono     & $+0.21$ \ $[+0.15,\,+0.27]$ & 0.57 & 3 / 0 / 11 & 1183 \\
\bottomrule
\end{tabular}
\end{table}

\begin{table}[H]
\centering
\caption{Diferencia de potencia media $\Delta_k$ [dB] (especular menos
monoestática) por posición en el Monte Carlo. $^{\ast}$: el intervalo de
confianza del 95\,\% excluye el cero.}
\label{tab:mc_posicion}
\begin{tabular}{@{}rrrrrr@{}}
\toprule
$\epsilon$ [°] & \textbf{Cubo} & \textbf{Caja} & \textbf{Cilindro} & \textbf{Esfera} & \textbf{Cono} \\
\midrule
20 & $-0.54$\sig & $-0.28$ & $-0.83$\sig & $-0.21$\sig & $+0.66$\sig \\
25 & $-0.74$\sig & $-0.25$ & $-0.89$\sig & $-0.17$\sig & $+0.85$\sig \\
30 & $-0.50$\sig & $-0.62$\sig & $-0.75$\sig & $-0.23$\sig & $+0.22$ \\
35 & $-0.47$\sig & $-0.34$ & $-0.62$\sig & $-0.45$\sig & $+0.26$\sig \\
40 & $-0.37$ & $-0.34$ & $-0.59$\sig & $-0.79$\sig & $-0.03$ \\
45 & $-0.35$ & $-0.30$ & $-0.42$\sig & $-1.17$\sig & $+0.19$ \\
50 & $-0.37$ & $-0.35$\sig & $-0.21$\sig & $-1.35$\sig & $+0.17$ \\
55 & $-0.04$ & $-0.10$ & $+0.03$ & $-1.21$\sig & $+0.03$ \\
60 & $-0.14$ & $-0.36$\sig & $+0.10$ & $-0.72$\sig & $+0.17$ \\
65 & $-0.21$ & $+0.14$ & $+0.17$\sig & $-0.15$\sig & $+0.08$ \\
70 & $-0.15$ & $+0.13$ & $+0.24$\sig & $+0.28$\sig & $-0.13$ \\
75 & $+0.33$\sig & $+0.03$ & $+0.21$\sig & $+0.42$\sig & $+0.06$ \\
80 & $+0.07$ & $+0.22$\sig & $+0.10$\sig & $+0.30$\sig & $+0.09$ \\
85 & $+0.03$ & $+0.08$ & $+0.03$ & $+0.10$\sig & $+0.02$ \\
\bottomrule
\end{tabular}
\end{table}

\begin{figure}[H]
\centering
\includegraphics[width=\textwidth]{fig_resumen_objetos.png}
\caption{Diferencia energética entre la configuración biestática especular
y la monoestática para los cinco objetos. (a) Objeto en orientación
nominal. (b) Monte Carlo con orientación aleatoria (media sobre las 14
posiciones fuera del nadir e intervalo de confianza del 95\,\%). Obsérvese
la diferencia de escala vertical entre ambos paneles. El desglose por
posición de ambos paneles se presenta en la Figura~\ref{fig:resumen_angulo}.}
\label{fig:resumen}
\end{figure}

La esfera merece un comentario: al ser invariante frente a rotaciones, sus
repeticiones solo difieren por la orientación de las facetas, de modo que
los intervalos son muy estrechos y todas las diferencias resultan
``significativas''. La variación sistemática de $\Delta_k$ con la posición
(de $-1.35$ a $+0.42\dB$) refleja la dependencia de la RCS biestática PO de
una esfera de $ka\approx9$ con el ángulo biestático, que en el límite de la
óptica geométrica sería nula ($\sigma=\pi a^2$ para cualquier ángulo
biestático alejado de la dispersión hacia delante \cite{knott1993}). En la
esfera, el criterio de acierto especular no es aplicable, porque toda
orientación presenta el mismo retorno.

\subsection{Desglose por posición y por banda de elevación}
\label{sec:res_mc_angulo}

La Figura~\ref{fig:resumen_angulo} reúne, para los cinco objetos y en la
misma escala de elevación, la diferencia por posición en orientación
nominal, la diferencia de medias del Monte Carlo con su intervalo de
confianza y la probabilidad de superioridad $\Pr\{P_r^B>P_r^M\}$. Para
resumir la tendencia angular sin perder la resolución estadística, la
Tabla~\ref{tab:mc_bandas} agrega la energía en tres bandas: incidencias
rasantes y entorno de Brewster ($\epsilon=20^\circ$--$35^\circ$,
$\theta_i=55^\circ$--$70^\circ$), incidencias intermedias
($40^\circ$--$60^\circ$) y próximas al nadir ($65^\circ$--$85^\circ$), con
intervalos de confianza por \emph{bootstrap} pareado calculados igual que
el del agregado global.

\begin{figure}[H]
\centering
\includegraphics[width=0.9\textwidth]{fig_resumen_por_angulo.png}
\caption{Desglose por posición de la comparación especular--monoestática
para los cinco objetos. (a) Diferencia de potencia en orientación nominal
(Tabla~\ref{tab:nominal_posicion}). (b) Diferencia de potencia media en el
Monte Carlo con su intervalo de confianza del 95\,\% (los marcadores se
desplazan ligeramente en abscisa para que no se solapen). (c) Fracción de
repeticiones en que el especular recibe más que el monoestático. La línea
discontinua marca $\epsilon_B$ y las punteadas, los límites de las bandas de
la Tabla~\ref{tab:mc_bandas}. Obsérvese la diferencia de escala vertical
entre (a) y (b).}
\label{fig:resumen_angulo}
\end{figure}

\begin{table}[H]
\centering
\caption{Diferencia energética especular menos monoestática por banda de
elevación. Nominal: objeto sin rotar. Monte Carlo: $10^4$ orientaciones por
posición, con IC\,95\,\% bootstrap pareado; $^{\ast}$: el intervalo excluye
el cero. $\Pr\{B>M\}$: fracción de repeticiones de la banda en que el
especular recibe más.}
\label{tab:mc_bandas}
\small
\begin{tabular}{@{}llrrrr@{}}
\toprule
\textbf{Objeto} & \textbf{Banda} $\bm{\epsilon}$ & \textbf{Nominal [dB]} & \textbf{Monte Carlo [dB]} & \textbf{IC 95\,\% [dB]} & $\bm{\Pr\{B>M\}}$ \\
\midrule
\multirow{3}{*}{Cubo}     & $20^\circ$--$35^\circ$ & $+18.17$ & $-0.56$\sig & $[-0.76,\,-0.37]$ & 0.52 \\
                          & $40^\circ$--$60^\circ$ & $+16.48$ & $-0.26$\sig & $[-0.44,\,-0.08]$ & 0.52 \\
                          & $65^\circ$--$85^\circ$ & $+5.58$  & $+0.01$     & $[-0.13,\,+0.15]$ & 0.51 \\
\midrule
\multirow{3}{*}{Caja}     & $20^\circ$--$35^\circ$ & $+16.74$ & $-0.37$\sig & $[-0.54,\,-0.22]$ & 0.51 \\
                          & $40^\circ$--$60^\circ$ & $+16.59$ & $-0.29$\sig & $[-0.44,\,-0.15]$ & 0.50 \\
                          & $65^\circ$--$85^\circ$ & $+5.61$  & $+0.12$\sig & $[+0.01,\,+0.23]$ & 0.50 \\
\midrule
\multirow{3}{*}{Cilindro} & $20^\circ$--$35^\circ$ & $+18.16$ & $-0.77$\sig & $[-0.86,\,-0.69]$ & 0.49 \\
                          & $40^\circ$--$60^\circ$ & $+13.63$ & $-0.22$\sig & $[-0.30,\,-0.16]$ & 0.49 \\
                          & $65^\circ$--$85^\circ$ & $+2.11$  & $+0.15$\sig & $[+0.11,\,+0.19]$ & 0.53 \\
\midrule
\multirow{3}{*}{Esfera}   & $20^\circ$--$35^\circ$ & $+0.31$  & $-0.26$\sig & $[-0.27,\,-0.26]$ & 0.29 \\
                          & $40^\circ$--$60^\circ$ & $-0.94$  & $-1.04$\sig & $[-1.05,\,-1.04]$ & 0.02 \\
                          & $65^\circ$--$85^\circ$ & $-1.07$  & $+0.19$\sig & $[+0.19,\,+0.20]$ & 0.81 \\
\midrule
\multirow{3}{*}{Cono}     & $20^\circ$--$35^\circ$ & $-6.19$  & $+0.50$\sig & $[+0.38,\,+0.62]$ & 0.59 \\
                          & $40^\circ$--$60^\circ$ & $+3.27$  & $+0.10$     & $[0.00,\,+0.21]$ & 0.58 \\
                          & $65^\circ$--$85^\circ$ & $+3.21$  & $+0.02$     & $[-0.04,\,+0.09]$ & 0.55 \\
\bottomrule
\end{tabular}
\end{table}

El desglose muestra que, aunque el agregado de cada objeto es casi nulo, la
diferencia no es uniforme a lo largo del barrido:
\begin{itemize}[itemsep=2pt]
  \item \textbf{Cubo, caja y cilindro} presentan la misma tendencia: el
  monoestático recibe algo más de energía en las incidencias rasantes y en
  el entorno de Brewster ($-0.37$ a $-0.77\dB$ en $20^\circ$--$35^\circ$) y
  la diferencia se anula o cambia de signo cerca del nadir ($+0.01$ a
  $+0.15\dB$ en $65^\circ$--$85^\circ$). La tendencia crece con el ángulo
  biestático $\beta$, lo que es consistente con que la compensación entre la reducción
  del pico especular y el ensanchamiento de su lóbulo
  (Sección~\ref{sec:disc_promedio}) no es completa para caras de pocas
  longitudes de onda. En todas las posiciones, sin embargo,
  $\Pr\{B>M\}$ se mantiene entre 0.48 y 0.54: las diferencias de media se
  deben a la cola de eventos especulares, no al caso típico, cuya mediana es
  igual o mayor en el especular (hasta $1.4\dB$ en las posiciones rasantes;
  Apéndice~\ref{ap:figuras}).
  \item \textbf{Esfera:} es el único objeto en que el modo preferente queda
  determinado por la posición, porque su RCS casi no depende de la
  orientación. El monoestático es mejor entre $20^\circ$ y $65^\circ$, con
  un máximo de $-1.35\dB$ en $50^\circ$, donde $\Pr\{B>M\}$ cae por debajo
  de 0.01, y el especular es mejor entre $70^\circ$ y $85^\circ$ (hasta
  $+0.42\dB$, $\Pr\{B>M\}\approx0.98$).
  \item \textbf{Cono:} al contrario que los objetos con caras planas, el
  especular es mejor en las incidencias rasantes ($+0.50\dB$ en
  $20^\circ$--$35^\circ$, con un máximo de $+0.85\dB$ en $\epsilon=25^\circ$),
  y la diferencia no es significativa en el resto. Es también el único
  objeto con $\Pr\{B>M\}$ sistemáticamente superior a 0.5 (0.51--0.60).
\end{itemize}
En ninguna posición ni banda la diferencia del Monte Carlo alcanza
$1.4\dB$ en valor absoluto, frente a los $+16$ a $+18\dB$ que el
especular obtiene en la banda rasante en orientación nominal para los
objetos con cara horizontal. El desglose,
por tanto, no altera la conclusión del agregado: con orientación
desconocida ninguna configuración ofrece una ventaja energética relevante
en ninguna parte del barrido. Sí indica que, si se quisiera elegir, la
geometría especular es preferible cerca del nadir y la monoestática en las
incidencias rasantes, justo donde se sitúa el ángulo de Brewster.

\section{Sensibilidad al rango de orientaciones del Monte Carlo}
\label{sec:sensibilidad}

Se repitió el experimento con giros
restringidos a $\pm45^\circ$, que cubren todas las orientaciones distintas
del cubo por su simetría de $90^\circ$, pero no las de los demás objetos.
Los resultados del cubo, la caja, el cilindro y la esfera fueron
equivalentes ($-0.23$, $-0.11$, $-0.20$ y $-0.35\dB$, respectivamente),
mientras que el cono dio $-6.6\dB$ sin ningún acierto especular en el modo
biestático. Este último valor es un artefacto del muestreo: con
$\pm45^\circ$ ninguna superficie del cono puede quedar horizontal (su base
mira siempre hacia abajo y su generatriz está inclinada $73^\circ$ respecto
de la horizontal), por lo que la reflexión especular hacia el receptor
simétrico es geométricamente imposible. Con el rango completo
($\pm180^\circ$) el resultado del cono se alinea con el de los demás
objetos. La conclusión es que la distribución de orientaciones debe
cubrir todo el espacio de orientaciones del objeto; restringirla introduce
un sesgo que puede favorecer artificialmente a cualquiera de los dos modos.

% ══════════════════════════════════════════════════════════════════════════
\chapter{Discusión}
\label{cap:discusion}

\section{Por qué la ventaja nominal desaparece al promediar la orientación}
\label{sec:disc_promedio}

La ventaja biestática observada en orientación nominal tiene un origen
puramente geométrico. En la configuración especular, la bisectriz entre las
direcciones hacia el transmisor y el receptor es siempre vertical, y por el
teorema de equivalencia monoestático--biestático \cite{kell1965,willis1991}
la RCS biestática se aproxima a la monoestática observada a lo largo de esa
bisectriz, reducida por un factor que depende del ángulo biestático. Una
cara horizontal se ve así en reflexión especular en toda la trayectoria, con
una pérdida de solo $20\log_{10}\cos\theta_t\approx-1.1\dB$ en el extremo
del barrido, mientras que el monoestático solo la ve en el nadir. La
ventaja es, en consecuencia, una propiedad de la \emph{alineación} entre la
geometría de adquisición y una cara del blanco, no de la configuración
biestática en sí.

Cuando la orientación del blanco es aleatoria, la condición de reflexión
especular exige en el caso monoestático que una normal de cara apunte hacia
el radar, y en el biestático que apunte a lo largo de la bisectriz. Ambos
eventos son igualmente probables bajo una distribución de orientaciones que
cubra todo el espacio. Además, el pico especular biestático de una cara es
menor que el monoestático por un factor $\cos^2(\beta/2)$, pero su lóbulo en
el espacio de orientaciones de la normal es más ancho, en proporción
inversa, porque el vector de fase de \eqref{eq:po} tiene módulo
$2k_s\cos(\beta/2)$ en lugar de $2k_s$. Las dos dependencias se compensan en
el promedio sobre orientaciones, de modo que la energía media esperada es
esencialmente la misma para ambos modos. Los resultados numéricos
(diferencias menores de $0.4\dB$ para los cinco objetos) confirman este
argumento para objetos de formas muy distintas, y coinciden con el límite de
óptica geométrica, en el que la RCS de una esfera es independiente del
ángulo biestático.

\section{El papel real del ángulo de Brewster}

La ganancia de transmitancia en el ángulo de Brewster respecto del nadir es
de solo $+1.02\dB$ con $n_2=2$, idéntica para ambos modos. Este valor es
pequeño frente a las variaciones de la RCS del blanco con la geometría
(decenas de dB) y frente a la incertidumbre sobre su orientación. En
cambio, el ángulo de Brewster es decisivo en la configuración especular por
otra vía: anula la reflexión especular de la superficie, que de otro modo
supera al eco del blanco en 8--17\,dB en casi todo el barrido. La
optimización de trayectoria guiada por Brewster debe interpretarse, por
tanto, como una estrategia de \emph{supresión del retorno de superficie} y
no principalmente como una estrategia de maximización de la energía
transmitida. Esta supresión, además, solo es efectiva en un intervalo
angular estrecho en torno a $\theta_B$, lo que impone requisitos de
precisión de trayectoria más estrictos que los de la propia reflexión
especular del blanco (del orden del grado frente a decenas de grados a
425\,MHz).

Este resultado matiza la motivación habitual de la literatura revisada. Los
trabajos de GPR biestático \cite{orfeo2019} y las observaciones de
amortiguamiento de Brewster en SAR de superficie \cite{ouchi2017,ouchi2018,thirion2018}
destacan el ángulo de incidencia como criterio para minimizar la reflexión
en la interfaz; los resultados de este relatorio muestran que, para blancos
enterrados observados con geometría especular, el beneficio principal es
efectivamente la reducción del retorno de superficie, mientras que la
ganancia en la energía que alcanza el blanco es marginal.

\section{Implicaciones para el diseño de trayectorias}

\begin{enumerate}[itemsep=3pt]
  \item Una función de costo que premie la potencia recibida con la RCS
  nominal de un blanco con caras planas horizontales sobrestima
  sistemáticamente la ventaja biestática. Para blancos de orientación
  desconocida, el término de potencia debería usar la RCS promediada sobre
  orientaciones, para la que ambos modos son equivalentes.
  \item El beneficio robusto de la geometría biestática especular es la
  supresión del retorno de superficie en Brewster. Si el objetivo es la
  detección en presencia de una superficie reflectante, conviene incluir
  explícitamente la SCR (o un término equivalente de \emph{clutter} de
  superficie) en la función de costo.
  \item La tolerancia de posicionamiento es asimétrica: el lóbulo especular
  del blanco admite desplazamientos del receptor de 50 a 87\,m sobre el arco,
  mientras que la anulación de la reflexión de superficie exige mantener el
  ángulo de incidencia a menos de un grado de $\theta_B$.
  \item Para blancos con geometría conocida y orientación controlada
  ---por ejemplo, tuberías o estructuras con techo horizontal---, la
  configuración especular sí aporta ganancias de 6 a 10\,dB, y su uso está
  justificado.
\end{enumerate}

\section{Limitaciones}

Las conclusiones están sujetas a las hipótesis de la
Sección~\ref{sec:limitaciones}. Las más relevantes para su alcance son: (i)
PO subestima o sobrestima en varios dB los retornos no especulares de
objetos de pocas longitudes de onda, como muestra la validación con MoM,
aunque la comparación entre modos se basa principalmente en las
contribuciones especulares, que PO describe bien; (ii) la interfaz lisa
representa el caso más favorable para la anulación de la reflexión de
superficie en Brewster, y la rugosidad dejará un residuo difuso; (iii) el
análisis es de onda continua, polarización TM y un único blanco; y (iv) la
distribución de orientaciones supuesta (giros independientes y uniformes en
dos ejes) no es uniforme en el grupo de rotaciones; la coincidencia de
resultados con distintos rangos para los objetos simétricos sugiere que la
conclusión no depende de este detalle, pero una distribución uniforme en
$SO(3)$ sería la elección canónica para trabajos futuros.

% ══════════════════════════════════════════════════════════════════════════
\chapter{Conclusiones y trabajo futuro}
\label{cap:conclusiones}

\section{Conclusiones}

\begin{enumerate}[itemsep=4pt]
  \item \textbf{Literatura.} De los 22 trabajos revisados, ninguno combina
  SAR biestático tomográfico, refracción en la interfaz aire--suelo y
  optimización de la geometría de adquisición guiada por el ángulo de
  Brewster; los antecedentes más cercanos son el GPR biestático de par fijo
  y las observaciones pasivas del amortiguamiento de Brewster en SAR
  monoestático de superficie.

  \item \textbf{Transmitancia.} Con $n_2=2$, el ángulo de Brewster
  ($\theta_B=63.43^\circ$) aporta una ganancia de transmitancia de ida y
  vuelta de $+1.02\dB$ respecto del nadir, idéntica en las configuraciones
  monoestática y biestática especular.

  \item \textbf{Orientación nominal.} Para blancos con una cara plana
  horizontal (cubo, caja, cilindro de eje vertical), la configuración
  especular recibe entre $+6.2$ y $+9.6\dB$ más energía que la monoestática,
  porque la bisectriz vertical
  mantiene la reflexión especular de esa cara en toda la trayectoria. Para
  la esfera y el cono la diferencia es pequeña o de signo variable.

  \item \textbf{Orientación aleatoria.} Promediando sobre $10^4$
  orientaciones aleatorias por posición, la diferencia energética entre
  ambas configuraciones es despreciable para los cinco objetos (entre
  $-0.37$ y $+0.21\dB$), y cada modo supera al otro en aproximadamente la
  mitad de las repeticiones. La ventaja biestática nominal es, por tanto,
  consecuencia de la alineación entre la geometría y el blanco, y no una
  propiedad intrínseca de la configuración especular.

  \item \textbf{Superficie.} El beneficio robusto de la geometría especular
  en el ángulo de Brewster es la anulación de la reflexión especular de la
  superficie, que en el resto del barrido supera al eco del blanco en
  8--17\,dB; la relación blanco--superficie mejora en más de $40\dB$ en el
  entorno inmediato de $\theta_B$.

  \item \textbf{Metodología.} La media de la potencia recibida con
  orientación aleatoria está dominada por eventos de reflexión especular
  poco frecuentes; su estimación exige un número de repeticiones
  inversamente proporcional a la probabilidad de esos eventos, y una
  distribución de orientaciones que cubra todo el espacio de orientaciones
  del objeto.
\end{enumerate}

\section{Trabajo futuro}

\begin{enumerate}[itemsep=3pt]
  \item Incorporar la relación blanco--superficie a la función de costo del
  optimizador de trayectoria y evaluar la precisión angular requerida para
  mantener la supresión en presencia de errores de navegación.
  \item Modelar la rugosidad de la superficie para cuantificar el residuo de
  reflexión en Brewster y la relación blanco--\emph{clutter} realista.
  \item Validar las RCS con métodos de onda completa (MoM) para los objetos
  y frecuencias de interés, y extender el análisis a polarización TE y
  cruzada.
  \item Repetir el Monte Carlo con una distribución de orientaciones
  uniforme en $SO(3)$ y con distribuciones informadas por el tipo de blanco
  (por ejemplo, tuberías de eje horizontal).
  \item Incluir pérdidas dieléctricas, el factor de divergencia por
  refracción y el análisis de banda ancha, y extender el estudio a la
  trayectoria helicoidal completa con formación de imagen tomográfica.
\end{enumerate}

% ══════════════════════════════════════════════════════════════════════════
\begin{thebibliography}{99}

\bibitem{andre2025} D.~Andre, V.~P.~Krishnan, C.~Nolan, ``Artefact Analysis of Multistatic Three-Dimensional SAR Imaging with Two Linear Trajectories,'' arXiv:2503.11323, 2025.

\bibitem{arikan1989} O.~Arikan, D.~C.~Munson~Jr., ``A tomographic formulation of bistatic synthetic aperture radar,'' in \emph{Advances in Communications and Signal Processing}, W.~A.~Porter, S.~C.~Kak, Eds., Lecture Notes in Control and Information Sciences, vol.~129, pp.~418--431, Springer, 1989.

\bibitem{chen2022} Y.~Chen, R.~Chen, H.~Liu, J.~Guo, Y.~Wang, J.~Zhang, ``Research on Configuration Constraints of Airborne Bistatic SARs,'' \emph{Sensors}, vol.~22, no.~17, art.~6534, 2022. doi:10.3390/s22176534

\bibitem{chenz2024} Z.~Chen, Y.~Li, C.~Hu, et~al., ``Repeat-pass space-surface bistatic SAR tomography: accurate imaging and first experiment,'' \emph{Science China Information Sciences}, vol.~67, art.~192304, 2024. doi:10.1007/s11432-024-4089-2

\bibitem{ding2002} Y.~Ding, D.~C.~Munson~Jr., ``A fast back-projection algorithm for bistatic SAR imaging,'' in \emph{Proc. IEEE Int. Conf. on Image Processing (ICIP)}, vol.~II, pp.~449--452, 2002.

\bibitem{efron1993} B.~Efron, R.~J.~Tibshirani, \emph{An Introduction to the Bootstrap}. New York: Chapman \& Hall, 1993.

\bibitem{garza2010} G.~Garza, Z.~Qiao, ``Resolution analysis of bistatic SAR,'' in \emph{Proc. SPIE 8021, Radar Sensor Technology XV}, 80210F, 2010.

\bibitem{ge2019} N.~Ge, X.~X.~Zhu, ``Bistatic-like Differential SAR Tomography,'' \emph{IEEE Transactions on Geoscience and Remote Sensing}, vol.~57, no.~8, pp.~5883--5893, 2019.

\bibitem{gordon1975} W.~B.~Gordon, ``Far-field approximations to the Kirchhoff--Helmholtz representations of scattered fields,'' \emph{IEEE Transactions on Antennas and Propagation}, vol.~23, no.~4, pp.~590--592, 1975.

\bibitem{grathwohl2023} A.~Grathwohl, B.~Meinecke, M.~Widmann, J.~Kanz, C.~Waldschmidt, ``UAV-Based Bistatic SAR-Imaging Using a Stationary Repeater,'' \emph{IEEE Journal of Microwaves}, vol.~3, no.~2, pp.~625--634, 2023. doi:10.1109/JMW.2023.3253667

\bibitem{he2025} J.~He, H.~Xie, H.~Liu, Z.~Wu, B.~Xu, N.~Zhu, Z.~Lu, P.~Qin, ``Multi-Baseline Bistatic SAR Three-Dimensional Imaging Method Based on Phase Error Calibration Combining PGA and EB-ISOA,'' \emph{Remote Sensing}, vol.~17, no.~3, art.~363, 2025. doi:10.3390/rs17030363

\bibitem{homer2005} J.~Homer, E.~Donskoi, B.~Mojarrabi, J.~Palmer, K.~Kubik, ``Three-dimensional bistatic synthetic aperture radar imaging system: spatial resolution analysis,'' \emph{IEE Proceedings -- Radar, Sonar and Navigation}, vol.~152, no.~6, pp.~391--394, 2005.

\bibitem{horne2002} A.~M.~Horne, G.~Yates, ``Bistatic synthetic aperture radar,'' in \emph{Proc. RADAR 2002}, Edinburgh, IEE Publication No.~490, pp.~6--10, 2002.

\bibitem{kanz2025} J.~Kanz, C.~Gesell, C.~Bonfert, D.~Werbunat, A.~Grathwohl, J.~Aguilar, M.~Vossiek, C.~Waldschmidt, ``UAV-Borne Digital Radar System for Coherent Multistatic SAR Imaging,'' arXiv:2507.20792, 2025.

\bibitem{kell1965} R.~E.~Kell, ``On the derivation of bistatic RCS from monostatic measurements,'' \emph{Proceedings of the IEEE}, vol.~53, no.~8, pp.~983--988, 1965.

\bibitem{knott1993} E.~F.~Knott, J.~F.~Shaeffer, M.~T.~Tuley, \emph{Radar Cross Section}, 2nd ed. Norwood, MA: Artech House, 1993.

\bibitem{liu2024} M.~Liu et~al., ``Unmanned Airborne Bistatic Interferometric Synthetic Aperture Radar Data Processing Method Using Bi-Directional Synchronization Chain Signals,'' \emph{Remote Sensing}, vol.~16, no.~5, art.~769, 2024. doi:10.3390/rs16050769

\bibitem{orfeo2019} D.~Orfeo, Y.~Zhang, D.~Burns, J.~S.~Miller, D.~R.~Huston, T.~Xia, ``Bistatic Antenna Configurations for Air-Launched Ground Penetrating Radar,'' \emph{Journal of Applied Remote Sensing}, vol.~13, no.~2, art.~027501, 2019. doi:10.1117/1.JRS.13.027501

\bibitem{ouchi2017} K.~Ouchi, C.-S.~Yang, ``Direct Evidence and Implications of Brewster's Angle Damping Observed by Synthetic Aperture Radar Images,'' in \emph{Proc. IGARSS 2017}, 2017.

\bibitem{ouchi2018} K.~Ouchi, C.-S.~Yang, ``Brewster Angle Damping Observed in the TerraSAR-X Synthetic Aperture Radar Images of Man-Made Targets,'' \emph{IEEE Geoscience and Remote Sensing Letters}, vol.~15, no.~4, pp.~532--536, 2018.

\bibitem{sun2022} H.~Sun, Z.~Sun, T.~Chen, Y.~Miao, J.~Wu, J.~Yang, ``An Efficient Backprojection Algorithm Based on Wavenumber-Domain Spectral Splicing for Monostatic and Bistatic SAR Configurations,'' \emph{Remote Sensing}, vol.~14, no.~8, art.~1885, 2022. doi:10.3390/rs14081885

\bibitem{thirion2018} L.~Thirion-Lefevre, R.~Guinvarc'h, ``The Double Brewster Angle Effect,'' \emph{Comptes Rendus Physique}, vol.~19, no.~1--2, pp.~43--53, 2018. doi:10.1016/j.crhy.2018.02.003

\bibitem{wang2020} Z.~Wang, M.~Liu, G.~Ai, P.~Wang, K.~Lv, ``Focusing of Bistatic SAR With Curved Trajectory Based on Extended Azimuth Nonlinear Chirp Scaling,'' \emph{IEEE Transactions on Geoscience and Remote Sensing}, vol.~58, no.~6, pp.~4160--4179, 2020.

\bibitem{willis1991} N.~J.~Willis, \emph{Bistatic Radar}. Norwood, MA: Artech House, 1991.

\bibitem{wu2025} J.~Wu, J.~Yang, Z.~Li, W.~Pu, Z.~Sun, H.~An, Q.~Yang, ``Review of Bistatic Synthetic Aperture Radar Imaging Methods,'' \emph{Journal of Radars}, vol.~14, no.~5, pp.~1115--1141, 2025. doi:10.12000/JR25067

\bibitem{zhang2018} H.~Zhang, J.~Tang, R.~Wang, Y.~Deng, W.~Wang, N.~Li, ``An Accelerated Backprojection Algorithm for Monostatic and Bistatic SAR Processing,'' \emph{Remote Sensing}, vol.~10, no.~1, art.~140, 2018. doi:10.3390/rs10010140

\bibitem{zhao2025} T.~Zhao, J.~Wang, Z.~Cheng, Z.~Huang, J.~Song, ``A Subimage Autofocus Bistatic Ground Cartesian Back-Projection Algorithm for Passive Bistatic SAR Based on GEO Satellites,'' \emph{Remote Sensing}, vol.~17, no.~9, art.~1576, 2025. doi:10.3390/rs17091576

\end{thebibliography}

% ══════════════════════════════════════════════════════════════════════════
\appendix
\chapter{Resultados complementarios por objeto}
\label{ap:figuras}

Este apéndice reúne, para cada objeto simulado, su geometría, los
diagramas 3D de RCS en el suelo, el balance en orientación nominal y el
resultado del Monte Carlo, estos dos últimos con su tabla de valores por
posición, que permite identificar en qué elevaciones es mejor cada
configuración. En los diagramas 3D, el radio representa
$\sigma$ en dB por encima de un piso situado 30\,dB por debajo del máximo
común a ambos diagramas del objeto; la flecha morada indica la dirección del
transmisor en Brewster, la naranja la dirección especular del receptor y el
cono punteado las direcciones accesibles desde el aire. Los paneles del
balance nominal y del Monte Carlo se describen en las
Figuras~\ref{fig:cubo_principal} y~\ref{fig:cubo_mc}.

\section{Cubo}
\label{ap:cubo}

\subsection{Geometría simulada}

Cubo PEC de 1\,m de arista con las caras alineadas con los ejes; las seis caras planas se describen exactamente con dos triángulos cada una. Su cara superior horizontal es la responsable de la reflexión especular analizada en el Capítulo~\ref{cap:resultados}.

La geometría del cubo se muestra en la Figura~\ref{fig:cubo_geometria} (Sección~\ref{sec:res_nominal}).

\subsection{Diagramas 3D de RCS}

\begin{figure}[H]
\centering
\begin{subfigure}{0.49\textwidth}\includegraphics[width=\textwidth]{fig_cubo_rcs3d_mono.png}\caption{Monoestático}\end{subfigure}\hfill
\begin{subfigure}{0.49\textwidth}\includegraphics[width=\textwidth]{fig_cubo_rcs3d_bi.png}\caption{Biestático, Tx en Brewster}\end{subfigure}
\caption{Diagramas 3D de RCS en el suelo (425\,MHz): cubo de 1\,m de arista.}
\label{fig:rcs3d_cubo}
\end{figure}

\subsection{Balance en orientación nominal}

El balance del cubo en orientación nominal se presenta en la Tabla~\ref{tab:cubo_nominal} y en la Figura~\ref{fig:cubo_principal} (Sección~\ref{sec:res_nominal}).

\subsection{Monte Carlo}

El resultado del Monte Carlo del cubo se presenta en la Figura~\ref{fig:cubo_mc} (Sección~\ref{sec:res_mc}).
El desglose por posición se da en la Tabla~\ref{tab:cubo_mc}.

\begin{table}[H]
\centering
\caption{Monte Carlo por posición ($10^4$ orientaciones): cubo de 1\,m. $\bar P_r$: potencia media; $\tilde P_r$: mediana; $\Delta_k$: diferencia de medias especular menos monoestática ($^{\ast}$: el IC\,95\,\% excluye el cero); aciertos: repeticiones con reflexión especular (M / B).}
\label{tab:cubo_mc}
\small
\setlength{\tabcolsep}{4pt}
\begin{tabular}{@{}rrrrcrrcc@{}}
\toprule
$\epsilon$ & $\bar P_r^M$ & $\bar P_r^B$ & $\Delta_k$ & IC 95\,\% & $\tilde P_r^M$ & $\tilde P_r^B$ & $\Pr\{B>M\}$ & Aciertos \\
{[°]} & [dBm] & [dBm] & [dB] & [dB] & [dBm] & [dBm] & & M / B \\
\midrule
20 & $-16.28$ & $-16.83$ & $-0.54$\sig & $[-0.93,\,-0.18]$ & $-38.37$ & $-37.00$ & 0.53 & 1031 / 1104 \\
25 & $-16.36$ & $-17.10$ & $-0.74$\sig & $[-1.12,\,-0.36]$ & $-38.38$ & $-37.11$ & 0.51 & 1015 / 1015 \\
30 & $-16.57$ & $-17.07$ & $-0.50$\sig & $[-0.89,\,-0.09]$ & $-38.52$ & $-37.13$ & 0.50 & 964 / 1029 \\
35 & $-16.57$ & $-17.04$ & $-0.47$\sig & $[-0.88,\,-0.07]$ & $-38.71$ & $-37.18$ & 0.52 & 983 / 1025 \\
40 & $-16.63$ & $-17.00$ & $-0.37$ & $[-0.75,\,+0.03]$ & $-38.68$ & $-37.22$ & 0.53 & 1008 / 1041 \\
45 & $-16.80$ & $-17.15$ & $-0.35$ & $[-0.75,\,+0.05]$ & $-38.82$ & $-37.46$ & 0.54 & 1004 / 1039 \\
50 & $-16.87$ & $-17.24$ & $-0.37$ & $[-0.76,\,+0.03]$ & $-38.96$ & $-37.62$ & 0.52 & 1017 / 1020 \\
55 & $-17.37$ & $-17.41$ & $-0.04$ & $[-0.47,\,+0.36]$ & $-39.55$ & $-38.55$ & 0.52 & 933 / 981 \\
60 & $-17.21$ & $-17.35$ & $-0.14$ & $[-0.54,\,+0.25]$ & $-39.03$ & $-38.41$ & 0.50 & 980 / 1020 \\
65 & $-17.29$ & $-17.51$ & $-0.21$ & $[-0.61,\,+0.18]$ & $-39.21$ & $-38.41$ & 0.52 & 964 / 987 \\
70 & $-17.40$ & $-17.55$ & $-0.15$ & $[-0.51,\,+0.22]$ & $-39.68$ & $-39.11$ & 0.53 & 980 / 1000 \\
75 & $-17.69$ & $-17.36$ & $+0.33$\sig & $[+0.00,\,+0.66]$ & $-39.47$ & $-39.04$ & 0.51 & 960 / 992 \\
80 & $-17.36$ & $-17.29$ & $+0.07$ & $[-0.17,\,+0.31]$ & $-39.55$ & $-39.79$ & 0.50 & 1008 / 1035 \\
85 & $-17.63$ & $-17.60$ & $+0.03$ & $[-0.10,\,+0.16]$ & $-39.57$ & $-39.42$ & 0.51 & 968 / 966 \\
\bottomrule
\end{tabular}
\end{table}

\clearpage
\section{Caja}
\label{ap:caja}

\subsection{Geometría simulada}

Paralelepípedo PEC de $1\times0.5$\,m en planta y 0.3\,m de altura. Su cara superior tiene la mitad del área de la del cubo, por lo que su retorno especular vertical es $6\dB$ menor ($14.0\dBsm$).

\begin{figure}[H]
\centering
\includegraphics[width=0.92\textwidth]{fig_caja_geometria.png}
\caption{Geometría simulada: caja de $1\times0.5\times0.3$\,m (12 triángulos). (a) Malla de facetas. (b) Sección en el plano $xz$ con el objeto enterrado.}
\label{fig:caja_geometria}
\end{figure}

\subsection{Diagramas 3D de RCS}

\begin{figure}[H]
\centering
\begin{subfigure}{0.49\textwidth}\includegraphics[width=\textwidth]{fig_caja_rcs3d_mono.png}\caption{Monoestático}\end{subfigure}\hfill
\begin{subfigure}{0.49\textwidth}\includegraphics[width=\textwidth]{fig_caja_rcs3d_bi.png}\caption{Biestático, Tx en Brewster}\end{subfigure}
\caption{Diagramas 3D de RCS en el suelo (425\,MHz): caja de $1\times0.5\times0.3$\,m.}
\label{fig:rcs3d_caja}
\end{figure}

\subsection{Balance en orientación nominal}

\begin{figure}[H]
\centering
\includegraphics[width=\textwidth]{fig_caja_principal.png}
\caption{Balance en orientación nominal: caja de $1\times0.5\times0.3$\,m.}
\label{fig:caja_principal}
\end{figure}

\begin{table}[H]
\centering
\caption{Balance por posición en orientación nominal: caja de $1\times0.5\times0.3$\,m. SCR: relación entre la potencia biestática del blanco y la reflexión especular de la superficie.}
\label{tab:caja_nominal}
\small
\begin{tabular}{@{}rrrrrrr@{}}
\toprule
$\epsilon$ & $\sigma^M$ & $\sigma^B$ & $P_r^M$ & $P_r^B$ & $\Delta_k$ & SCR \\
{[°]} & [dBsm] & [dBsm] & [dBm] & [dBm] & [dB] & [dB] \\
\midrule
20 & $-4.40$ & $+12.94$ & $-28.76$ & $-11.41$ & $+17.35$ & $-14.11$ \\
25 & $-3.26$ & $+13.04$ & $-27.49$ & $-11.19$ & $+16.30$ & $+0.89$ \\
30 & $-2.96$ & $+13.15$ & $-27.23$ & $-11.12$ & $+16.11$ & $-7.00$ \\
35 & $-4.13$ & $+13.27$ & $-28.51$ & $-11.11$ & $+17.40$ & $-13.46$ \\
40 & $-8.45$ & $+13.38$ & $-32.96$ & $-11.13$ & $+21.83$ & $-16.52$ \\
45 & $-19.32$ & $+13.50$ & $-43.98$ & $-11.16$ & $+32.82$ & $-18.40$ \\
50 & $-5.26$ & $+13.61$ & $-30.06$ & $-11.19$ & $+18.87$ & $-19.68$ \\
55 & $-0.30$ & $+13.71$ & $-25.22$ & $-11.22$ & $+14.00$ & $-20.60$ \\
60 & $+0.84$ & $+13.79$ & $-24.20$ & $-11.25$ & $+12.95$ & $-21.28$ \\
65 & $-2.65$ & $+13.87$ & $-27.79$ & $-11.27$ & $+16.52$ & $-21.79$ \\
70 & $-12.71$ & $+13.93$ & $-37.93$ & $-11.29$ & $+26.64$ & $-22.16$ \\
75 & $+4.52$ & $+13.97$ & $-20.77$ & $-11.31$ & $+9.45$ & $-22.44$ \\
80 & $+10.32$ & $+14.01$ & $-15.01$ & $-11.33$ & $+3.68$ & $-22.62$ \\
85 & $+13.14$ & $+14.02$ & $-12.22$ & $-11.34$ & $+0.88$ & $-22.74$ \\
90 & $+14.02$ & $+14.02$ & $-11.35$ & $-11.35$ & $+0.00$ & $-22.78$ \\
\bottomrule
\end{tabular}
\end{table}

\subsection{Monte Carlo}

\begin{figure}[H]
\centering
\includegraphics[width=\textwidth]{fig_caja_montecarlo.png}
\caption{Monte Carlo con $10^4$ orientaciones aleatorias por posición: caja de $1\times0.5\times0.3$\,m.}
\label{fig:caja_mc}
\end{figure}

\begin{table}[H]
\centering
\caption{Monte Carlo por posición ($10^4$ orientaciones): caja de $1\times0.5\times0.3$\,m. $\bar P_r$: potencia media; $\tilde P_r$: mediana; $\Delta_k$: diferencia de medias especular menos monoestática ($^{\ast}$: el IC\,95\,\% excluye el cero); aciertos: repeticiones con reflexión especular (M / B).}
\label{tab:caja_mc}
\small
\setlength{\tabcolsep}{4pt}
\begin{tabular}{@{}rrrrcrrcc@{}}
\toprule
$\epsilon$ & $\bar P_r^M$ & $\bar P_r^B$ & $\Delta_k$ & IC 95\,\% & $\tilde P_r^M$ & $\tilde P_r^B$ & $\Pr\{B>M\}$ & Aciertos \\
{[°]} & [dBm] & [dBm] & [dB] & [dB] & [dBm] & [dBm] & & M / B \\
\midrule
20 & $-26.83$ & $-27.11$ & $-0.28$ & $[-0.60,\,+0.07]$ & $-36.86$ & $-36.17$ & 0.52 & 237 / 272 \\
25 & $-26.80$ & $-27.05$ & $-0.25$ & $[-0.59,\,+0.09]$ & $-36.83$ & $-36.31$ & 0.51 & 223 / 287 \\
30 & $-26.73$ & $-27.35$ & $-0.62$\sig & $[-0.94,\,-0.29]$ & $-36.93$ & $-36.55$ & 0.51 & 254 / 236 \\
35 & $-26.90$ & $-27.25$ & $-0.34$ & $[-0.69,\,+0.00]$ & $-37.07$ & $-36.56$ & 0.51 & 228 / 247 \\
40 & $-26.94$ & $-27.28$ & $-0.34$ & $[-0.68,\,+0.02]$ & $-37.25$ & $-36.81$ & 0.51 & 239 / 256 \\
45 & $-27.24$ & $-27.53$ & $-0.30$ & $[-0.60,\,+0.02]$ & $-37.13$ & $-36.89$ & 0.51 & 236 / 255 \\
50 & $-27.28$ & $-27.63$ & $-0.35$\sig & $[-0.68,\,-0.01]$ & $-37.51$ & $-37.47$ & 0.50 & 227 / 239 \\
55 & $-27.65$ & $-27.75$ & $-0.10$ & $[-0.42,\,+0.22]$ & $-37.99$ & $-38.01$ & 0.49 & 214 / 242 \\
60 & $-27.40$ & $-27.76$ & $-0.36$\sig & $[-0.66,\,-0.02]$ & $-37.58$ & $-37.51$ & 0.48 & 247 / 232 \\
65 & $-27.80$ & $-27.66$ & $+0.14$ & $[-0.16,\,+0.43]$ & $-37.61$ & $-37.63$ & 0.49 & 239 / 270 \\
70 & $-27.81$ & $-27.68$ & $+0.13$ & $[-0.16,\,+0.41]$ & $-37.74$ & $-37.63$ & 0.51 & 224 / 248 \\
75 & $-27.68$ & $-27.65$ & $+0.03$ & $[-0.24,\,+0.29]$ & $-37.87$ & $-37.78$ & 0.50 & 242 / 255 \\
80 & $-27.94$ & $-27.72$ & $+0.22$\sig & $[+0.03,\,+0.41]$ & $-37.99$ & $-38.09$ & 0.50 & 221 / 235 \\
85 & $-27.83$ & $-27.75$ & $+0.08$ & $[-0.04,\,+0.20]$ & $-38.03$ & $-37.87$ & 0.51 & 246 / 249 \\
\bottomrule
\end{tabular}
\end{table}

\clearpage
\section{Cilindro}
\label{ap:cilindro}

\subsection{Geometría simulada}

Cilindro PEC de eje vertical con tapas planas, radio 0.3\,m y longitud 1\,m. La superficie lateral se faceta con 21 segmentos por vuelta (flecha $<\lambda_s/100$). La tapa superior, un disco de $0.28\,\mathrm{m}^2$, actúa como cara especular horizontal.

\begin{figure}[H]
\centering
\includegraphics[width=0.92\textwidth]{fig_cilindro_geometria.png}
\caption{Geometría simulada: cilindro de eje vertical, $a=0.3$\,m, $L=1$\,m (84 triángulos). (a) Malla de facetas. (b) Sección en el plano $xz$ con el objeto enterrado.}
\label{fig:cilindro_geometria}
\end{figure}

\subsection{Diagramas 3D de RCS}

\begin{figure}[H]
\centering
\begin{subfigure}{0.49\textwidth}\includegraphics[width=\textwidth]{fig_cilindro_rcs3d_mono.png}\caption{Monoestático}\end{subfigure}\hfill
\begin{subfigure}{0.49\textwidth}\includegraphics[width=\textwidth]{fig_cilindro_rcs3d_bi.png}\caption{Biestático, Tx en Brewster}\end{subfigure}
\caption{Diagramas 3D de RCS en el suelo (425\,MHz): cilindro de eje vertical, $a=0.3$\,m, $L=1$\,m.}
\label{fig:rcs3d_cilindro}
\end{figure}

\subsection{Balance en orientación nominal}

\begin{figure}[H]
\centering
\includegraphics[width=\textwidth]{fig_cilindro_principal.png}
\caption{Balance en orientación nominal: cilindro de eje vertical, $a=0.3$\,m, $L=1$\,m.}
\label{fig:cilindro_principal}
\end{figure}

\begin{table}[H]
\centering
\caption{Balance por posición en orientación nominal: cilindro de eje vertical, $a=0.3$\,m, $L=1$\,m. SCR: relación entre la potencia biestática del blanco y la reflexión especular de la superficie.}
\label{tab:cilindro_nominal}
\small
\begin{tabular}{@{}rrrrrrr@{}}
\toprule
$\epsilon$ & $\sigma^M$ & $\sigma^B$ & $P_r^M$ & $P_r^B$ & $\Delta_k$ & SCR \\
{[°]} & [dBsm] & [dBsm] & [dBm] & [dBm] & [dB] & [dB] \\
\midrule
20 & $-9.81$ & $+7.85$ & $-34.17$ & $-16.50$ & $+17.66$ & $-19.21$ \\
25 & $-9.82$ & $+7.84$ & $-34.06$ & $-16.39$ & $+17.67$ & $-4.31$ \\
30 & $-10.20$ & $+7.86$ & $-34.46$ & $-16.41$ & $+18.05$ & $-12.29$ \\
35 & $-11.60$ & $+7.92$ & $-35.98$ & $-16.46$ & $+19.52$ & $-18.81$ \\
40 & $-15.76$ & $+8.02$ & $-40.28$ & $-16.49$ & $+23.79$ & $-21.89$ \\
45 & $-23.35$ & $+8.17$ & $-48.00$ & $-16.49$ & $+31.52$ & $-23.73$ \\
50 & $-10.48$ & $+8.34$ & $-35.28$ & $-16.46$ & $+18.82$ & $-24.95$ \\
55 & $-4.01$ & $+8.51$ & $-28.93$ & $-16.42$ & $+12.51$ & $-25.80$ \\
60 & $+0.13$ & $+8.66$ & $-24.91$ & $-16.39$ & $+8.53$ & $-26.42$ \\
65 & $+3.07$ & $+8.77$ & $-22.07$ & $-16.37$ & $+5.70$ & $-26.88$ \\
70 & $+5.26$ & $+8.86$ & $-19.97$ & $-16.37$ & $+3.60$ & $-27.23$ \\
75 & $+6.87$ & $+8.91$ & $-18.41$ & $-16.38$ & $+2.04$ & $-27.50$ \\
80 & $+8.01$ & $+8.94$ & $-17.32$ & $-16.39$ & $+0.93$ & $-27.69$ \\
85 & $+8.69$ & $+8.95$ & $-16.67$ & $-16.41$ & $+0.25$ & $-27.81$ \\
90 & $+8.94$ & $+8.94$ & $-16.43$ & $-16.43$ & $+0.00$ & $-27.86$ \\
\bottomrule
\end{tabular}
\end{table}

\subsection{Monte Carlo}

\begin{figure}[H]
\centering
\includegraphics[width=\textwidth]{fig_cilindro_montecarlo.png}
\caption{Monte Carlo con $10^4$ orientaciones aleatorias por posición: cilindro de eje vertical, $a=0.3$\,m, $L=1$\,m.}
\label{fig:cilindro_mc}
\end{figure}

\begin{table}[H]
\centering
\caption{Monte Carlo por posición ($10^4$ orientaciones): cilindro de eje vertical, $a=0.3$\,m, $L=1$\,m. $\bar P_r$: potencia media; $\tilde P_r$: mediana; $\Delta_k$: diferencia de medias especular menos monoestática ($^{\ast}$: el IC\,95\,\% excluye el cero); aciertos: repeticiones con reflexión especular (M / B).}
\label{tab:cilindro_mc}
\small
\setlength{\tabcolsep}{4pt}
\begin{tabular}{@{}rrrrcrrcc@{}}
\toprule
$\epsilon$ & $\bar P_r^M$ & $\bar P_r^B$ & $\Delta_k$ & IC 95\,\% & $\tilde P_r^M$ & $\tilde P_r^B$ & $\Pr\{B>M\}$ & Aciertos \\
{[°]} & [dBm] & [dBm] & [dB] & [dB] & [dBm] & [dBm] & & M / B \\
\midrule
20 & $-24.29$ & $-25.12$ & $-0.83$\sig & $[-1.01,\,-0.66]$ & $-34.68$ & $-34.00$ & 0.49 & 2568 / 2628 \\
25 & $-24.32$ & $-25.21$ & $-0.89$\sig & $[-1.07,\,-0.72]$ & $-34.74$ & $-33.86$ & 0.50 & 2490 / 2531 \\
30 & $-24.51$ & $-25.27$ & $-0.75$\sig & $[-0.93,\,-0.59]$ & $-34.66$ & $-34.03$ & 0.48 & 2408 / 2491 \\
35 & $-24.57$ & $-25.19$ & $-0.62$\sig & $[-0.80,\,-0.44]$ & $-34.82$ & $-34.04$ & 0.49 & 2414 / 2502 \\
40 & $-24.59$ & $-25.19$ & $-0.59$\sig & $[-0.75,\,-0.42]$ & $-34.79$ & $-34.27$ & 0.48 & 2464 / 2529 \\
45 & $-24.69$ & $-25.11$ & $-0.42$\sig & $[-0.57,\,-0.26]$ & $-35.06$ & $-34.37$ & 0.48 & 2484 / 2530 \\
50 & $-24.89$ & $-25.10$ & $-0.21$\sig & $[-0.36,\,-0.06]$ & $-35.18$ & $-34.55$ & 0.48 & 2478 / 2523 \\
55 & $-25.24$ & $-25.21$ & $+0.03$ & $[-0.14,\,+0.18]$ & $-35.44$ & $-34.98$ & 0.49 & 2330 / 2417 \\
60 & $-25.04$ & $-24.94$ & $+0.10$ & $[-0.04,\,+0.24]$ & $-35.34$ & $-35.00$ & 0.50 & 2530 / 2561 \\
65 & $-25.20$ & $-25.02$ & $+0.17$\sig & $[+0.05,\,+0.31]$ & $-35.46$ & $-35.16$ & 0.53 & 2481 / 2536 \\
70 & $-25.36$ & $-25.12$ & $+0.24$\sig & $[+0.12,\,+0.35]$ & $-35.54$ & $-35.42$ & 0.54 & 2448 / 2527 \\
75 & $-25.47$ & $-25.26$ & $+0.21$\sig & $[+0.11,\,+0.31]$ & $-35.69$ & $-35.56$ & 0.54 & 2419 / 2444 \\
80 & $-25.45$ & $-25.36$ & $+0.10$\sig & $[+0.02,\,+0.17]$ & $-35.71$ & $-35.69$ & 0.52 & 2452 / 2437 \\
85 & $-25.44$ & $-25.40$ & $+0.03$ & $[0.00,\,+0.07]$ & $-35.69$ & $-35.64$ & 0.51 & 2454 / 2480 \\
\bottomrule
\end{tabular}
\end{table}

\clearpage
\section{Esfera}
\label{ap:esfera}

\subsection{Geometría simulada}

Esfera PEC de 0.5\,m de radio, facetada en 27 meridianos y 14 paralelos (flecha $<\lambda_s/100$). Al ser invariante frente a rotaciones, sirve de referencia sin orientación preferente.

\begin{figure}[H]
\centering
\includegraphics[width=0.92\textwidth]{fig_esfera_geometria.png}
\caption{Geometría simulada: esfera de radio $a=0.5$\,m (702 triángulos). (a) Malla de facetas. (b) Sección en el plano $xz$ con el objeto enterrado.}
\label{fig:esfera_geometria}
\end{figure}

\subsection{Diagramas 3D de RCS}

\begin{figure}[H]
\centering
\begin{subfigure}{0.49\textwidth}\includegraphics[width=\textwidth]{fig_esfera_rcs3d_mono.png}\caption{Monoestático}\end{subfigure}\hfill
\begin{subfigure}{0.49\textwidth}\includegraphics[width=\textwidth]{fig_esfera_rcs3d_bi.png}\caption{Biestático, Tx en Brewster}\end{subfigure}
\caption{Diagramas 3D de RCS en el suelo (425\,MHz): esfera de radio $a=0.5$\,m.}
\label{fig:rcs3d_esfera}
\end{figure}

\subsection{Balance en orientación nominal}

\begin{figure}[H]
\centering
\includegraphics[width=\textwidth]{fig_esfera_principal.png}
\caption{Balance en orientación nominal: esfera de radio $a=0.5$\,m.}
\label{fig:esfera_principal}
\end{figure}

\begin{table}[H]
\centering
\caption{Balance por posición en orientación nominal: esfera de radio $a=0.5$\,m. SCR: relación entre la potencia biestática del blanco y la reflexión especular de la superficie.}
\label{tab:esfera_nominal}
\small
\begin{tabular}{@{}rrrrrrr@{}}
\toprule
$\epsilon$ & $\sigma^M$ & $\sigma^B$ & $P_r^M$ & $P_r^B$ & $\Delta_k$ & SCR \\
{[°]} & [dBsm] & [dBsm] & [dBm] & [dBm] & [dB] & [dB] \\
\midrule
20 & $-1.06$ & $-1.21$ & $-25.41$ & $-25.56$ & $-0.15$ & $-28.26$ \\
25 & $-1.06$ & $-0.89$ & $-25.29$ & $-25.13$ & $+0.17$ & $-13.04$ \\
30 & $-1.03$ & $-0.14$ & $-25.29$ & $-24.41$ & $+0.89$ & $-20.29$ \\
35 & $-0.94$ & $-0.68$ & $-25.31$ & $-25.05$ & $+0.26$ & $-27.41$ \\
40 & $-0.84$ & $-0.35$ & $-25.35$ & $-24.86$ & $+0.49$ & $-30.26$ \\
45 & $-0.79$ & $-1.16$ & $-25.44$ & $-25.81$ & $-0.37$ & $-33.06$ \\
50 & $-0.88$ & $-2.91$ & $-25.68$ & $-27.70$ & $-2.03$ & $-36.20$ \\
55 & $-1.05$ & $-3.02$ & $-25.98$ & $-27.94$ & $-1.97$ & $-37.32$ \\
60 & $-1.12$ & $-2.88$ & $-26.17$ & $-27.92$ & $-1.75$ & $-37.95$ \\
65 & $-1.02$ & $-2.67$ & $-26.16$ & $-27.81$ & $-1.65$ & $-38.33$ \\
70 & $-0.81$ & $-2.50$ & $-26.03$ & $-27.72$ & $-1.69$ & $-38.59$ \\
75 & $-0.69$ & $-1.97$ & $-25.98$ & $-27.25$ & $-1.27$ & $-38.38$ \\
80 & $-0.97$ & $-1.55$ & $-26.30$ & $-26.88$ & $-0.58$ & $-38.18$ \\
85 & $-1.37$ & $-1.55$ & $-26.73$ & $-26.91$ & $-0.18$ & $-38.30$ \\
90 & $-1.54$ & $-1.54$ & $-26.91$ & $-26.91$ & $+0.00$ & $-38.34$ \\
\bottomrule
\end{tabular}
\end{table}

\subsection{Monte Carlo}

\begin{figure}[H]
\centering
\includegraphics[width=\textwidth]{fig_esfera_montecarlo.png}
\caption{Monte Carlo con $10^4$ orientaciones aleatorias por posición: esfera de radio $a=0.5$\,m.}
\label{fig:esfera_mc}
\end{figure}

\begin{table}[H]
\centering
\caption{Monte Carlo por posición ($10^4$ orientaciones): esfera de radio $a=0.5$\,m. $\bar P_r$: potencia media; $\tilde P_r$: mediana; $\Delta_k$: diferencia de medias especular menos monoestática ($^{\ast}$: el IC\,95\,\% excluye el cero); aciertos: repeticiones con reflexión especular (M / B).}
\label{tab:esfera_mc}
\small
\setlength{\tabcolsep}{4pt}
\begin{tabular}{@{}rrrrcrrcc@{}}
\toprule
$\epsilon$ & $\bar P_r^M$ & $\bar P_r^B$ & $\Delta_k$ & IC 95\,\% & $\tilde P_r^M$ & $\tilde P_r^B$ & $\Pr\{B>M\}$ & Aciertos \\
{[°]} & [dBm] & [dBm] & [dB] & [dB] & [dBm] & [dBm] & & M / B \\
\midrule
20 & $-25.20$ & $-25.41$ & $-0.21$\sig & $[-0.22,\,-0.19]$ & $-25.19$ & $-25.43$ & 0.33 & 10000 / 10000 \\
25 & $-25.09$ & $-25.26$ & $-0.17$\sig & $[-0.18,\,-0.16]$ & $-25.08$ & $-25.28$ & 0.35 & 10000 / 10000 \\
30 & $-25.12$ & $-25.35$ & $-0.23$\sig & $[-0.24,\,-0.22]$ & $-25.11$ & $-25.36$ & 0.30 & 10000 / 10000 \\
35 & $-25.23$ & $-25.68$ & $-0.45$\sig & $[-0.46,\,-0.44]$ & $-25.22$ & $-25.69$ & 0.16 & 10000 / 10000 \\
40 & $-25.36$ & $-26.15$ & $-0.79$\sig & $[-0.80,\,-0.77]$ & $-25.35$ & $-26.15$ & 0.06 & 10000 / 10000 \\
45 & $-25.51$ & $-26.67$ & $-1.17$\sig & $[-1.18,\,-1.16]$ & $-25.50$ & $-26.68$ & 0.02 & 10000 / 10000 \\
50 & $-25.65$ & $-27.00$ & $-1.35$\sig & $[-1.36,\,-1.34]$ & $-25.64$ & $-27.02$ & 0.00 & 10000 / 10000 \\
55 & $-25.78$ & $-26.99$ & $-1.21$\sig & $[-1.22,\,-1.20]$ & $-25.77$ & $-27.00$ & 0.00 & 10000 / 10000 \\
60 & $-25.89$ & $-26.62$ & $-0.72$\sig & $[-0.73,\,-0.72]$ & $-25.88$ & $-26.61$ & 0.01 & 10000 / 10000 \\
65 & $-25.99$ & $-26.14$ & $-0.15$\sig & $[-0.16,\,-0.14]$ & $-25.98$ & $-26.13$ & 0.28 & 10000 / 10000 \\
70 & $-26.07$ & $-25.80$ & $+0.28$\sig & $[+0.27,\,+0.28]$ & $-26.06$ & $-25.77$ & 0.90 & 10000 / 10000 \\
75 & $-26.14$ & $-25.72$ & $+0.42$\sig & $[+0.42,\,+0.43]$ & $-26.13$ & $-25.70$ & 0.98 & 10000 / 10000 \\
80 & $-26.18$ & $-25.88$ & $+0.30$\sig & $[+0.30,\,+0.31]$ & $-26.17$ & $-25.87$ & 0.98 & 10000 / 10000 \\
85 & $-26.21$ & $-26.11$ & $+0.10$\sig & $[+0.10,\,+0.10]$ & $-26.20$ & $-26.10$ & 0.91 & 10000 / 10000 \\
\bottomrule
\end{tabular}
\end{table}

\clearpage
\section{Cono}
\label{ap:cono}

\subsection{Geometría simulada}

Cono PEC de base circular de 0.3\,m de radio y 1\,m de altura, con el vértice hacia arriba. No tiene ninguna cara horizontal orientada hacia la superficie: su generatriz forma $73^\circ$ con la horizontal y su base mira hacia abajo.

\begin{figure}[H]
\centering
\includegraphics[width=0.92\textwidth]{fig_cono_geometria.png}
\caption{Geometría simulada: cono de base $a=0.3$\,m y altura $L=1$\,m, vértice hacia arriba (42 triángulos). (a) Malla de facetas. (b) Sección en el plano $xz$ con el objeto enterrado.}
\label{fig:cono_geometria}
\end{figure}

\subsection{Diagramas 3D de RCS}

\begin{figure}[H]
\centering
\begin{subfigure}{0.49\textwidth}\includegraphics[width=\textwidth]{fig_cono_rcs3d_mono.png}\caption{Monoestático}\end{subfigure}\hfill
\begin{subfigure}{0.49\textwidth}\includegraphics[width=\textwidth]{fig_cono_rcs3d_bi.png}\caption{Biestático, Tx en Brewster}\end{subfigure}
\caption{Diagramas 3D de RCS en el suelo (425\,MHz): cono de base $a=0.3$\,m y altura $L=1$\,m, vértice hacia arriba.}
\label{fig:rcs3d_cono}
\end{figure}

\subsection{Balance en orientación nominal}

\begin{figure}[H]
\centering
\includegraphics[width=\textwidth]{fig_cono_principal.png}
\caption{Balance en orientación nominal: cono de base $a=0.3$\,m y altura $L=1$\,m, vértice hacia arriba.}
\label{fig:cono_principal}
\end{figure}

\begin{table}[H]
\centering
\caption{Balance por posición en orientación nominal: cono de base $a=0.3$\,m y altura $L=1$\,m, vértice hacia arriba. SCR: relación entre la potencia biestática del blanco y la reflexión especular de la superficie.}
\label{tab:cono_nominal}
\small
\begin{tabular}{@{}rrrrrrr@{}}
\toprule
$\epsilon$ & $\sigma^M$ & $\sigma^B$ & $P_r^M$ & $P_r^B$ & $\Delta_k$ & SCR \\
{[°]} & [dBsm] & [dBsm] & [dBm] & [dBm] & [dB] & [dB] \\
\midrule
20 & $-21.41$ & $-31.42$ & $-45.76$ & $-55.77$ & $-10.01$ & $-58.47$ \\
25 & $-21.37$ & $-30.42$ & $-45.61$ & $-54.66$ & $-9.05$ & $-42.57$ \\
30 & $-21.27$ & $-29.40$ & $-45.54$ & $-53.66$ & $-8.12$ & $-49.55$ \\
35 & $-21.14$ & $-23.55$ & $-45.51$ & $-47.93$ & $-2.41$ & $-50.28$ \\
40 & $-21.20$ & $-23.09$ & $-45.71$ & $-47.60$ & $-1.89$ & $-52.99$ \\
45 & $-21.57$ & $-22.57$ & $-46.22$ & $-47.22$ & $-1.00$ & $-54.47$ \\
50 & $-22.16$ & $-19.35$ & $-46.95$ & $-44.15$ & $+2.80$ & $-52.64$ \\
55 & $-22.85$ & $-17.04$ & $-47.77$ & $-41.96$ & $+5.81$ & $-51.34$ \\
60 & $-23.51$ & $-16.08$ & $-48.55$ & $-41.12$ & $+7.43$ & $-51.15$ \\
65 & $-23.65$ & $-16.00$ & $-48.79$ & $-41.14$ & $+7.65$ & $-51.65$ \\
70 & $-22.17$ & $-15.94$ & $-47.40$ & $-41.16$ & $+6.24$ & $-52.03$ \\
75 & $-19.73$ & $-15.90$ & $-45.02$ & $-41.18$ & $+3.84$ & $-52.31$ \\
80 & $-17.62$ & $-15.87$ & $-42.95$ & $-41.21$ & $+1.75$ & $-52.50$ \\
85 & $-16.30$ & $-15.86$ & $-41.66$ & $-41.22$ & $+0.44$ & $-52.62$ \\
90 & $-15.86$ & $-15.86$ & $-41.23$ & $-41.23$ & $+0.00$ & $-52.65$ \\
\bottomrule
\end{tabular}
\end{table}

\subsection{Monte Carlo}

\begin{figure}[H]
\centering
\includegraphics[width=\textwidth]{fig_cono_montecarlo.png}
\caption{Monte Carlo con $10^4$ orientaciones aleatorias por posición: cono de base $a=0.3$\,m y altura $L=1$\,m, vértice hacia arriba.}
\label{fig:cono_mc}
\end{figure}

\begin{table}[H]
\centering
\caption{Monte Carlo por posición ($10^4$ orientaciones): cono de base $a=0.3$\,m y altura $L=1$\,m, vértice hacia arriba. $\bar P_r$: potencia media; $\tilde P_r$: mediana; $\Delta_k$: diferencia de medias especular menos monoestática ($^{\ast}$: el IC\,95\,\% excluye el cero); aciertos: repeticiones con reflexión especular (M / B).}
\label{tab:cono_mc}
\small
\setlength{\tabcolsep}{4pt}
\begin{tabular}{@{}rrrrcrrcc@{}}
\toprule
$\epsilon$ & $\bar P_r^M$ & $\bar P_r^B$ & $\Delta_k$ & IC 95\,\% & $\tilde P_r^M$ & $\tilde P_r^B$ & $\Pr\{B>M\}$ & Aciertos \\
{[°]} & [dBm] & [dBm] & [dB] & [dB] & [dBm] & [dBm] & & M / B \\
\midrule
20 & $-29.55$ & $-28.89$ & $+0.66$\sig & $[+0.41,\,+0.91]$ & $-39.86$ & $-37.50$ & 0.59 & 1224 / 1487 \\
25 & $-29.57$ & $-28.72$ & $+0.85$\sig & $[+0.60,\,+1.10]$ & $-39.93$ & $-37.45$ & 0.60 & 1183 / 1467 \\
30 & $-29.38$ & $-29.16$ & $+0.22$ & $[-0.02,\,+0.47]$ & $-39.94$ & $-38.15$ & 0.58 & 1236 / 1351 \\
35 & $-29.26$ & $-29.00$ & $+0.26$\sig & $[+0.01,\,+0.50]$ & $-39.82$ & $-38.00$ & 0.58 & 1327 / 1461 \\
40 & $-29.42$ & $-29.45$ & $-0.03$ & $[-0.28,\,+0.20]$ & $-40.00$ & $-38.37$ & 0.59 & 1287 / 1342 \\
45 & $-29.61$ & $-29.42$ & $+0.19$ & $[-0.05,\,+0.43]$ & $-40.00$ & $-38.43$ & 0.59 & 1302 / 1415 \\
50 & $-29.83$ & $-29.66$ & $+0.17$ & $[-0.07,\,+0.41]$ & $-40.29$ & $-38.86$ & 0.58 & 1244 / 1349 \\
55 & $-29.88$ & $-29.85$ & $+0.03$ & $[-0.21,\,+0.25]$ & $-40.42$ & $-39.50$ & 0.56 & 1328 / 1322 \\
60 & $-30.15$ & $-29.97$ & $+0.17$ & $[-0.04,\,+0.39]$ & $-40.55$ & $-39.76$ & 0.57 & 1232 / 1301 \\
65 & $-30.08$ & $-29.99$ & $+0.08$ & $[-0.12,\,+0.27]$ & $-40.64$ & $-39.99$ & 0.56 & 1311 / 1325 \\
70 & $-30.17$ & $-30.30$ & $-0.13$ & $[-0.30,\,+0.06]$ & $-40.76$ & $-40.37$ & 0.56 & 1302 / 1260 \\
75 & $-30.37$ & $-30.30$ & $+0.06$ & $[-0.09,\,+0.21]$ & $-40.94$ & $-40.58$ & 0.56 & 1242 / 1245 \\
80 & $-30.27$ & $-30.18$ & $+0.09$ & $[-0.01,\,+0.19]$ & $-40.72$ & $-40.55$ & 0.54 & 1296 / 1320 \\
85 & $-30.41$ & $-30.39$ & $+0.02$ & $[-0.04,\,+0.08]$ & $-40.92$ & $-40.80$ & 0.51 & 1270 / 1272 \\
\bottomrule
\end{tabular}
\end{table}

\end{document}
